% ===========================================================================
%  5 (continuação) — Modelagem UML
%
%  Sem \chapter DE PROPÓSITO (ADR-122): este arquivo continua o Capítulo 5.
%  Um capítulo novo depois do 5 renumeraria do 6 ao 21, e "cap. 17" é citado
%  dezenas de vezes em comentário de código da observabilidade --- renumerar
%  não quebra a citação, faz ela apontar para outro capítulo em silêncio.
%
%  Todo diagrama daqui foi conferido contra o código em 2026-09-27 e cita os
%  arquivos de onde saiu. Mudou a peça? O diagrama muda no MESMO commit, e o
%  par Mermaid no Apêndice A também.
% ===========================================================================
\section{Modelagem UML}
\label{sec:uml}

O C4 das seções anteriores responde \emph{que peças existem e quem fala com
quem}. Ele não tem notação --- por escolha de quem o desenhou~\cite{c4model} ---
para o que uma peça \emph{é} (tipo, multiplicidade, herança, interface), para o
que ela \emph{faz ao longo do tempo} (estado, prazo, ordem de mensagens) nem para
\emph{uma situação concreta} (instâncias com valores). É o que a UML cobre. Esta
seção desenha a plataforma nos catorze tipos de diagrama da UML
2.5.1~\cite{omg2017uml}, distribuídos pelas cinco visões do modelo
4+1~\cite{kruchten1995views} que abrem este capítulo (\adrr{ADR-122}).

Nem todo tipo rende o mesmo, e a seção não finge que rende. Onde o documento já
tinha um diagrama daquele tipo --- casos de uso, sequência, máquina de estados,
atividades ---, a Tabela~\ref{tab:uml-catalogo} aponta para ele e a seção
acrescenta só o que faltava. Onde não tinha, cada diagrama foi posto no lugar em
que o tipo mostra algo que nenhum outro mostra. O de tempo, por exemplo, é o
único em que a janela de 401 entre o \emph{token} de 15 minutos e o pedido de
renovação a cada 10 aparece sem conta de cabeça (Figura~\ref{fig:uml-tempo}).

\begin{regra}
\textbf{Diagrama que diverge do código é pior do que diagrama nenhum}, porque
quem lê acredita nele. Cada figura desta seção cita os arquivos de onde saiu, e
todas foram conferidas contra o código em 2026-09-27. Quem muda a peça muda o
diagrama e o par Mermaid do Apêndice~\ref{ap:mermaid} no mesmo \emph{commit}. Os
nomes são os do código (\cd{QuotaService}, \cd{exigirMissaoLiberada}), porque é
por eles que se procura; o texto em volta é em português.
\end{regra}

\subsection{Os catorze diagramas e onde cada um está}
\label{sec:uml-catalogo}

A UML separa os diagramas em dois ramos. Os \textbf{de estrutura} descrevem o que
existe e como se compõe, sem tempo. Os \textbf{de comportamento} descrevem o que
acontece, e quatro deles --- os \textbf{de interação} --- descrevem o que acontece
\emph{entre} participantes, por mensagens. A Figura~\ref{fig:uml-taxonomia} é a
própria taxonomia desenhada como diagrama de classes, com o número das figuras
deste documento em cada folha.

\begin{figure}[H]
\centering
\fitwidth{%
\begin{tikzpicture}[x=1cm, y=1cm]
  \tikzset{
    tx/.style={rectangle, draw=unooNavy, line width=0.8pt, fill=unooSurface,
      font=\scriptsize, text width=38mm, minimum height=9.8mm, inner sep=3pt,
      align=left},
    txabs/.style={tx, fill=unooNavy, text=white, align=center,
      font=\scriptsize\bfseries\itshape},
  }
  \def\fg#1{{\color{unooInkSoft}\scriptsize #1}}

  % --- raiz e os dois ramos
  \node[txabs] (diag) at (8.3,0) {Diagrama};
  \node[txabs] (est)  at (3.4,-1.9) {Diagrama de estrutura};
  \node[txabs] (comp) at (12.4,-1.9) {Diagrama de comportamento};
  \draw[associa] (est.north) -- (3.4,-0.95) -- (12.4,-0.95) -- (comp.north);
  \draw[hered] (8.3,-0.95) -- (diag.south);

  % --- estrutura: sete folhas num tronco
  \foreach \nome/\figs/\y in {
      Classes/{\ref{fig:uml-classes-dominio}, \ref{fig:uml-classes-ia}; antes: \ref{fig:c4l4}, \ref{fig:graus}}/-3.3,
      Objetos/{\ref{fig:uml-objetos}}/-4.6,
      Pacotes/{\ref{fig:uml-pacotes}; antes: \ref{fig:monorepo} (grafo)}/-5.9,
      Componentes/{\ref{fig:uml-componentes}}/-7.2,
      Estrutura composta/{\ref{fig:uml-estrutura}}/-8.5,
      Implantação/{\ref{fig:uml-implantacao}; antes: \ref{fig:deploy-contrato} (infra)}/-9.8,
      Perfil/{\ref{fig:uml-perfil}}/-11.1} {
    \node[tx, anchor=west] at (2.7,\y) {\textbf{\nome}\\\fg{Fig. \figs}};
    \draw[associa] (2.1,\y) -- (2.7,\y);
  }
  \draw[hered] (2.1,-11.1) -- (2.1,-2.39);

  % --- comportamento: tres folhas e o ramo abstrato da interacao
  \foreach \nome/\figs/\y in {
      Casos de uso/{\ref{fig:uml-casos-gestao}; antes: \ref{fig:casos-de-uso}}/-3.3,
      Atividades/{\ref{fig:uml-atividade-disp}; antes: \ref{fig:fluxo-embed}, \ref{fig:guardrails-bemestar}, \ref{fig:app92-caminho}}/-4.6,
      Máquina de estados/{\ref{fig:uml-estados-ponte}; antes: \ref{fig:estados-geracao}}/-5.9} {
    \node[tx, anchor=west] at (11.4,\y) {\textbf{\nome}\\\fg{Fig. \figs}};
    \draw[associa] (10.8,\y) -- (11.4,\y);
  }
  \node[txabs, anchor=west] (inter) at (11.4,-7.2) {Diagrama de interação};
  \draw[associa] (10.8,-7.2) -- (inter.west);
  \draw[hered] (10.8,-7.2) -- (10.8,-2.39);

  \foreach \nome/\figs/\y in {
      Sequência/{\ref{fig:uml-seq-entrada}, \ref{fig:uml-seq-ia}, \ref{fig:uml-seq-agente}, \ref{fig:uml-seq-ajuda}; antes: \ref{fig:seq-feed}, \ref{fig:seq-handshake}, \ref{fig:seq-erro-cliente}}/-8.75,
      Comunicação/{\ref{fig:uml-comunicacao}}/-10.15,
      Tempo/{\ref{fig:uml-tempo}}/-11.45,
      Visão geral de interação/{\ref{fig:uml-visao-geral}}/-12.75} {
    \node[tx, anchor=west, text width=34mm] at (12.6,\y) {\textbf{\nome}\\\fg{Fig. \figs}};
    \draw[associa] (12.0,\y) -- (12.6,\y);
  }
  \draw[hered] (12.0,-12.75) -- (12.0,-7.69);

  \node[actnote, text width=62mm, anchor=north west,
        execute at begin node={\hyphenpenalty=10000\exhyphenpenalty=10000}]
    at (2.1,-12.0)
    {\textbf{Como ler:} caixa escura em itálico é classificador
     \emph{abstrato} --- não existe diagrama \enquote{de estrutura} que não seja de
     um tipo concreto. \enquote{Antes} aponta para a figura do mesmo tipo que o
     documento já tinha; \enquote{(grafo)} e \enquote{(infra)} marcam as que
     cumprem o papel com outra notação.};
\end{tikzpicture}}
\caption[Os catorze tipos de diagrama da UML 2.5.1]{Os catorze tipos de diagrama
  da UML 2.5.1, desenhados como a especificação os organiza: uma hierarquia de
  generalização com três classificadores abstratos. Cada folha traz o número das
  figuras deste documento daquele tipo.}
\label{fig:uml-taxonomia}
\end{figure}

A Tabela~\ref{tab:uml-catalogo} diz, para cada tipo, a pergunta que ele responde
\emph{nesta} plataforma --- que é o critério pelo qual cada exemplar foi escolhido.

\begin{tabularx}{\linewidth}{@{}L{26mm}L{21mm}X L{16mm}@{}}
\caption{Os catorze diagramas UML e a pergunta que cada um responde aqui}
\label{tab:uml-catalogo}\\
\toprule
\headrow \thd{Diagrama} & \thd{Visão 4+1} & \thd{A pergunta que ele responde aqui} & \thd{Figuras}\\
\midrule
\endfirsthead
\toprule
\headrow \thd{Diagrama} & \thd{Visão 4+1} & \thd{A pergunta que ele responde aqui} & \thd{Figuras}\\
\midrule
\endhead
\multicolumn{4}{@{}l}{\textbf{\textcolor{unooNavy}{Estrutura}}}\\
Perfil & --- & Que estereótipos a plataforma acrescenta à UML, e que
  restrição cada um carrega? & \ref{fig:uml-perfil}\\
\zebra Classes & Lógica & Que conceitos existem, com que multiplicidade --- e onde a
  regra é restrição do modelo e não só do código? A porta da IA: que interface,
  que implementações, quem decora quem? & \ref{fig:uml-classes-dominio},
  \ref{fig:uml-classes-ia}\\
Objetos & Lógica & Numa rede concreta, com três linhas em
  \cd{disponibilidades}, o que fica ligado e por quê? & \ref{fig:uml-objetos}\\
\zebra Pacotes & \hspace{0pt}Desenvolvimento & Quem pode importar quem no \emph{monorepo} e dentro
  da API? & \ref{fig:uml-pacotes}\\
Componentes & \hspace{0pt}Desenvolvimento & Que interface cada peça \emph{fornece} e qual
  \emph{requer} --- quais delas são HTTP e qual é \cd{postMessage}? &
  \ref{fig:uml-componentes}\\
\zebra Estrutura composta & \hspace{0pt}Desenvolvimento & Por dentro da página da missão: que
  partes, que portas, e por onde o \emph{token} atravessa o \emph{sandbox}? &
  \ref{fig:uml-estrutura}\\
Implantação & Física & Que artefato roda em que nó, e por que caminho de
  comunicação? & \ref{fig:uml-implantacao}\\
\multicolumn{4}{@{}l}{\textbf{\textcolor{unooNavy}{Comportamento}}}\\
\zebra Casos de uso & Cenários & Quem administra a rede, e o que cada papel da
  cadeia pode fazer que o de baixo não pode? & \ref{fig:uml-casos-gestao},
  \ref{fig:casos-de-uso}\\
Atividades & Processos & Qual é o algoritmo que decide se um recurso está ligado,
  e que parte dele é pura? & \ref{fig:uml-atividade-disp}\\
\zebra Máquina de estados & Processos & Em que estados a ponte pode estar, do lado
  da experiência, e o que a tira de cada um? & \ref{fig:uml-estados-ponte},
  \ref{fig:estados-geracao}\\
Sequência & Processos & Em que ordem as duas travas da entrada na missão
  disparam? E a quota da IA: o que roda antes e o que roda depois do
  \emph{provider}? E quando quem chama o modelo é o \emph{vendor} do avatar? &
  \ref{fig:uml-seq-entrada}, \ref{fig:uml-seq-ia}, \ref{fig:uml-seq-agente},
  \ref{fig:uml-seq-ajuda}\\
\zebra Comunicação & Processos & Que objetos um lote de telemetria atravessa, em que
  ordem numerada, até virar linha no banco? & \ref{fig:uml-comunicacao}\\
Tempo & Processos & Quando o \emph{token}, a sessão e o relógio da SDK vencem, um em
  relação ao outro? & \ref{fig:uml-tempo}\\
\zebra Visão geral de interação & Processos & Como as interações anteriores se
  encadeiam numa jornada do estudante, do \emph{login} à missão? &
  \ref{fig:uml-visao-geral}\\
\bottomrule
\end{tabularx}

\subsection{O perfil da plataforma}
\label{sec:uml-perfil}

Um \textbf{perfil} é o mecanismo que a UML oferece para uma plataforma estender a
linguagem sem alterá-la: cada \emph{estereótipo} se aplica a uma metaclasse da UML
(\cd{Class}, \cd{Component}\ldots) e carrega \emph{valores etiquetados} e
\emph{restrições}. Ele vem primeiro porque os diagramas seguintes usam os
estereótipos dele, e porque escrever as regras invioláveis do \arq{AGENTS.md}
como restrição de estereótipo mostra a que elemento cada uma se aplica. As
restrições são as do código: \cd{escola\_id NOT NULL} e RLS forçada nas tabelas
com dado de estudante (Seção~\ref{sec:c4l4} e Capítulo~\ref{cap:tenancy}), a
guarda de liberação no funil das rotas de estudante (\adrr{ADR-69}), a quota antes
e depois da IA (\adrr{ADR-105}) e o caminho imutável do \emph{bundle}.

\begin{figure}[H]
\centering
\fitwidth{%
\begin{tikzpicture}[x=1cm, y=1cm]

  % sem hifenizacao: as caixas sao estreitas, a quebra vai a mao
  \hyphenpenalty=10000 \exhyphenpenalty=10000
  % --- estereotipos, em coluna; cada um logo abaixo do anterior
  \node[classe, text width=50mm, anchor=north west] (sTab) at (4.6,-1.2)
    {\ctitulo{estereótipo}{TabelaDeTenant}
     \nodepart{two}
     \texttt{prefixo: String [0..1]}\\
     {\color{unooMuted}código da experiência;\\ vazio na plataforma}
     \nodepart{three}
     \{\texttt{escola\_id} NOT NULL\}\\
     \{ENABLE e FORCE\\\ \ ROW LEVEL SECURITY\}\\
     \{toda \emph{query} dentro de \texttt{comTenancy}\}};

  \node[classe, text width=50mm, anchor=north west] (sExp)
      at ([yshift=-7mm]sTab.south west)
    {\ctitulo{estereótipo}{Experiencia}
     \nodepart{two}
     \texttt{id: String} {\color{unooMuted}(o \texttt{expId})}\\
     \texttt{versao: SemVer}\\
     \texttt{produto: Produto}\\
     \texttt{publico: Publico = estudante}\\
     \texttt{permissoes: String [*]}\\
     \texttt{viewport: String [0..1]}
     \nodepart{three}
     \{gabarito e correção só na API\}\\
     \{\emph{token} só em memória\}};

  \node[classe, text width=50mm, anchor=north west] (sBun)
      at ([yshift=-9mm]sExp.south west)
    {\ctitulo{estereótipo}{Bundle}
     \nodepart{two}
     \texttt{caminho = gs://\{contrato\}-exp/}\\
     \texttt{\ \ \{expId\}/\{versao\}/}
     \nodepart{three}
     \{imutável: mudar o conteúdo é subir a \texttt{versao}\}\\
     \{\texttt{versao} igual ao \texttt{version} do \texttt{package.json}\}};

  \node[classe, text width=50mm, anchor=north west] (sRota)
      at ([yshift=-7mm]sBun.south west)
    {\ctitulo{estereótipo}{RotaDeEstudante}
     \nodepart{two}
     \texttt{EXP\_ID: String}
     \nodepart{three}
     \{chama \texttt{exigirMissaoLiberada} no funil\}\\
     \{403 sem linha em \texttt{liberacoes}\}};

  \node[classe, text width=50mm, anchor=north west] (sIA)
      at ([xshift=8mm]sRota.north east)
    {\ctitulo{estereótipo}{ChamadaDeIA}
     \nodepart{two}
     \texttt{perfil: estudante | professor}
     \nodepart{three}
     \{\texttt{QuotaService.temSaldo} antes\}\\
     \{\texttt{QuotaService.consumir} depois\}\\
     \{vendor só atrás de uma porta de \texttt{ai/}:\\\ \ \texttt{AI\_PROVIDER} ou \texttt{AGENTE\_PROFESSOR}\}};

  % --- metaclasses da UML, fora do perfil, na altura do seu estereotipo
  \node[metaclasse] (mClass) at (sTab.west -| 1.6,0) {\estereo{metaclasse}\\\textbf{Class}};
  \node[metaclasse] (mComp)  at (sExp.west -| 1.6,0) {\estereo{metaclasse}\\\textbf{Component}};
  \node[metaclasse] (mArt)   at (sBun.west -| 1.6,0) {\estereo{metaclasse}\\\textbf{Artifact}};
  \node[metaclasse] (mOp)    at (sRota.west -| 1.6,0) {\estereo{metaclasse}\\\textbf{Operation}};

  % --- coluna da direita: nota, enumeracoes, nota
  \node[actnote, text width=50mm, anchor=north west] (n1) at (sTab.north east -| sIA.west)
    {\textbf{Como ler:} um elemento marcado \estereo{TabelaDeTenant} em outro
     diagrama carrega as três restrições sem repeti-las. A seta cheia é
     \emph{extensão} de metaclasse; não confundir com a herança, de seta
     aberta.};
  \node[enumeracao, text width=26mm, anchor=north west] (eProd) at (sExp.north east -| sIA.west)
    {\ctitulo{enumeração}{Produto}
     \nodepart{two}
     \texttt{entre-redes}\\
     \texttt{plano-de-aula}\\
     \texttt{plano-de-ensino}\\
     \texttt{bem-estar}\\
     \texttt{agente-professor}\\
     \texttt{ajuda}};
  \node[enumeracao, text width=20mm, anchor=north west] (ePub) at ([xshift=4mm]eProd.north east)
    {\ctitulo{enumeração}{Publico}
     \nodepart{two}
     \texttt{estudante}\\
     \texttt{professor}};
  \node[actnote, text width=50mm, anchor=north west] (n2) at (sBun.north east -| sIA.west)
    {\textbf{Onde o perfil ainda não é lei:} \texttt{geracao/} e as reações de
     \texttt{motor/} chamam o \emph{provider} sem \texttt{temSaldo} --- as
     exceções registradas na \textbf{ADR-74} e em \texttt{ai/AGENTS.md}. Uma
     \estereo{ChamadaDeIA} nova não herda a exceção.};

  % --- o pacote do perfil, no fundo
  \begin{scope}[on background layer]
    \pacote[fill=unooBlue!3]{perfil}{($(sTab.north west)+(-0.45,0.45)$)}%
      {($(sIA.south east)+(0.45,-0.95)$)}{\estereo{perfil} Unooria}
  \end{scope}

  % --- extensoes (seta cheia: estende a metaclasse)
  \draw[estende] (sTab.west) -- (mClass.east);
  \draw[estende] (sExp.west) -- (mComp.east);
  \draw[estende] (sBun.west) -- (mArt.east);
  \draw[estende] (sRota.west) -- (mOp.east);
  \draw[estende] (sIA.south) -- ++(0,-0.6) -| (mOp.south);

  % --- o artefato manifesta o componente; a experiencia usa as enumeracoes
  \draw[depend] (sBun.north) -- (sExp.south)
    node[cxlabel, midway, right=3pt]{\estereo{manifesta}};
  \draw[depend] (sExp.east |- eProd.west) -- (eProd.west);

\end{tikzpicture}}
\caption[Perfil UML da plataforma]{O perfil UML da plataforma: cinco estereótipos
  sobre quatro metaclasses. As restrições entre chaves são as regras
  invioláveis do \arq{AGENTS.md} ditas no lugar a que se aplicam. A Ajuda com
  avatar não é \estereo{ChamadaDeIA}: é a exceção única e deliberada à quota de
  \emph{tokens}, limitada por minutos de vídeo (\adrr{ADR-132}), e não um
  precedente para a próxima rota de IA. Fontes:
  \arq{recursos/catalogo.ts} (produtos), \arq{comum/liberacao.ts},
  \arq{ai/quota.service.ts}, \arq{prisma/prisma.service.ts} e
  \arq{public/experiencia.json} de cada experiência.}
\label{fig:uml-perfil}
\end{figure}

% ---------------------------------------------------------------------------
\subsection{Visão de cenários: a administração da rede}
\label{sec:uml-cenarios}

O diagrama de casos de uso do Capítulo~\ref{cap:casosdeuso}
(Figura~\ref{fig:casos-de-uso}) cobre a sala de aula e foi desenhado antes do
modelo de seis papéis (\adrr{ADR-80}). A Figura~\ref{fig:uml-casos-gestao} cobre
o que ele não cobre: quem administra a rede. A notação que carrega a regra é a
\textbf{generalização de atores}. Um ator filho faz tudo o que o pai faz, e o
código confirma a cadeia inteira: toda rota \cd{GESTORES} abre para os quatro
papéis de gestão, toda rota \cd{GESTORES\_DO\_CONTRATO} inclui o \cd{admin}, e
as únicas rotas que separam um papel do de cima são \cd{@Roles("admin")} (o painel
da Unooria, \adrr{ADR-88}) e o relatório da diretoria (\cd{gestor\_regional} em
diante, \adrr{ADR-82}). O que muda de um papel para o outro, fora essas rotas, é
o \emph{escopo} e não a capacidade: o \cd{gestor\_regional} trabalha uma escola
da regional por vez, e o \cd{gestor\_contrato} escolhe a escola por
\cd{?escolaId=} porque o \emph{token} dele não tem escola.

A cadeia de \emph{atribuição} é outra relação e não se confunde com a
generalização. Ela diz quem cria quem (\cd{ATRIBUI\_DIRETAMENTE}, em
\arq{usuarios/regras.ts}), e entra no diagrama como o caso de suporte
\ucr{UC-30}, incluído pelos dois caminhos de cadastro. O \ucr{UC-29} (balcão) não
escreve em \cd{users}: monta uma linha de planilha e inclui o \ucr{UC-20}, o
importador (\adrr{ADR-54}, \adrr{ADR-59}).

\begin{figure}[H]
\centering
\fitpage{%
\begin{tikzpicture}[x=1cm, y=1cm]

  % sem hifenizacao: a quebra das elipses vai a mao
  \hyphenpenalty=10000 \exhyphenpenalty=10000
  \tikzset{
    uc/.style={ucase, text width=, inner xsep=1.5pt, inner ysep=1pt,
      minimum width=36mm, minimum height=10mm},
    ucs/.style={uc, fill=unooSurface, draw=unooAxis},
    ucl/.style={ucnote, fill=unooSurface, inner sep=1.2pt},
  }

  % ---- fronteira do sistema
  \draw[ucsystem] (3.3,0.45) rectangle (13.75,-18.5);
  \node[font=\small\bfseries\color{unooNavy}, anchor=north] at (8.5,0.3)
    {Plataforma Unooria --- administração};

  % ---- atores da cadeia (esquerda): o filho faz tudo o que o pai faz
  \ator{gesc}{(1.5,-2.35)}{Gestor de escola}
  \ator{greg}{(1.5,-6.7)}{Gestor regional}
  \ator{gctr}{(1.5,-11.55)}{Gestor do contrato}
  \ator{gadm}{(1.5,-16.4)}{Admin (Unooria)}
  \draw[ucgen] (greg.west) -- ++(-0.95,0) |- (gesc.west);
  \draw[ucgen] (gctr.west) -- ++(-0.95,0) |- ([yshift=-2mm]greg.west);
  \draw[ucgen] (gadm.west) -- ++(-0.95,0) |- ([yshift=-2mm]gctr.west);

  % ---- atores da direita
  \ator{prof}{(15.3,-1.55)}{Professor}
  \atorsys{ci}{(15.3,-12.9)}{CI/CD}
  \ator{visit}{(15.3,-17.4)}{Visitante do site}

  % ---- coluna da esquerda: os casos da cadeia de gestao
  \node[uc] (uc29) at (5.8,-0.9)  {\ucr{UC-29} Cadastrar\\pessoa no balcão};
  \node[uc] (uc20) at (5.8,-2.5)  {\ucr{UC-20} Importar\\pessoas e enturmar};
  \node[uc] (uc21) at (5.8,-3.8)  {\ucr{UC-21} Editar pessoa\\e redefinir senha};
  \node[uc] (uc32) at (5.8,-5.4)  {\ucr{UC-32} Trocar\\a escola ativa};
  \node[uc] (uc33) at (5.8,-6.7)  {\ucr{UC-33} Ver o panorama\\da diretoria};
  \node[uc] (uc34) at (5.8,-8.0)  {\ucr{UC-34} Gerar o relatório\\da diretoria};
  \node[uc] (uc31) at (5.8,-9.6)  {\ucr{UC-31} Criar papel\\sem escola};
  \node[uc] (uc35) at (5.8,-10.9) {\ucr{UC-35} Manter escolas\\e regionais};
  \node[uc] (uc22) at (5.8,-12.2) {\ucr{UC-22} Ler o agregado\\da rede};
  \node[uc] (uc36) at (5.8,-13.5) {\ucr{UC-36} Ler o histórico\\de recursos};
  \node[uc] (uc37) at (5.8,-15.1) {\ucr{UC-37} Ligar ou\\desligar recurso};
  \node[uc] (uc39) at (5.8,-16.4) {\ucr{UC-39} Editar\\a quota de IA};
  \node[uc] (uc40) at (5.8,-17.7) {\ucr{UC-40} Ler cadastros\\de interesse};

  % ---- coluna da direita: professor, suporte, CI e visitante
  \node[uc]  (uc44) at (11.2,-0.9)  {\ucr{UC-44} Montar as\\equipes da missão};
  \node[uc]  (uc42) at (11.2,-2.2)  {\ucr{UC-42} Liberar ou encerrar\\missão na turma};
  \node[ucs] (uc43) at (11.2,-3.8)  {\ucr{UC-43} Conferir\\a disponibilidade};
  \node[ucs] (uc30) at (11.2,-6.7)  {\ucr{UC-30} Conferir a\\cadeia de atribuição};
  \node[uc]  (uc23) at (11.2,-12.9) {\ucr{UC-23} Publicar\\manifesto};
  \node[ucs] (uc24) at (11.2,-14.4) {\ucr{UC-24} Ativar ou\\reverter versão};
  \node[ucs] (uc38) at (11.2,-15.8) {\ucr{UC-38} Fixar versão\\por escopo};
  \node[uc]  (uc41) at (11.2,-17.4) {\ucr{UC-41} Cadastrar\\interesse};

  % ---- associacoes
  \foreach \u in {uc29, uc20, uc21} \draw[ucrel] (gesc) -- (\u.west);
  \foreach \u in {uc32, uc33, uc34} \draw[ucrel] (greg) -- (\u.west);
  \foreach \u in {uc31, uc35, uc22, uc36} \draw[ucrel] (gctr) -- (\u.west);
  \foreach \u in {uc37, uc39, uc40} \draw[ucrel] (gadm) -- (\u.west);
  \draw[ucrel] (prof) -- (uc44.east);
  \draw[ucrel] (prof) -- (uc42.east);
  \draw[ucrel] (ci) -- (uc23.east);
  \draw[ucrel] (visit) -- (uc41.east);

  % ---- include / extend
  \draw[ucinclude] (uc29.south) -- (uc20.north)
    node[ucl, midway, right=3pt]{\estereo{include}};
  \draw[ucinclude] (uc29.east) -- (8.55,-0.9) |- (uc30.west)
    node[ucl, pos=0.3]{\estereo{include}};
  \draw[ucinclude] (uc31.east) -| (uc30.south)
    node[ucl, pos=0.78, right=3pt]{\estereo{include}};
  \draw[ucinclude] (uc42.south) -- (uc43.north)
    node[ucl, midway, right=3pt]{\estereo{include}};
  \draw[ucextend] (uc24.north) -- (uc23.south)
    node[ucl, midway, right=3pt]{\estereo{extend}};
  \draw[ucextend] (uc24.west) -- (uc37.12)
    node[ucl, pos=0.45, above=2pt]{\estereo{extend}};
  \draw[ucextend] (uc38.west) -- (uc37.-12)
    node[ucl, pos=0.45, below=2pt]{\estereo{extend}};

\end{tikzpicture}}
\caption[Casos de uso da administração da rede]{Casos de uso da administração da
  rede, com os seis papéis da \adrr{ADR-80}. A seta de generalização entre atores é
  a regra: o \emph{admin} faz tudo o que o gestor do contrato faz, e mais o painel
  da Unooria (\ucr{UC-37} a \ucr{UC-40}, \adrr{ADR-88}); o gestor do contrato faz
  tudo o que o regional faz, com a rede inteira como escopo. Fontes:
  \arq{auth/claims.ts} (\cd{GESTORES}, \cd{GESTORES\_DO\_CONTRATO}),
  \arq{usuarios/regras.ts} (\cd{ATRIBUI\_DIRETAMENTE}) e os \cd{@Roles} de cada
  \emph{controller}.}
\label{fig:uml-casos-gestao}
\end{figure}

\begin{legendadiag}
\amostra{ucase} caso de uso primário \quad
\amostra{ucasesec} caso de uso de suporte (incluído ou extensão) \quad
seta de triângulo aberto entre atores: generalização
\end{legendadiag}

\begin{nota}
\textbf{Três ausências deliberadas.} O professor não herda de gestor nenhum: a
gestão dele é da turma em que leciona (\ucr{UC-42} exige matrícula de professor
ativa, e o gestor só lê a sala). \enquote{Ver o consumo de IA} não existe como
caso de uso: a tela de \ucr{UC-39} mostra os tetos, e o contador
\cd{ia\_uso\_diario} não sai por rota nenhuma. E o \ucr{UC-24} estende dois
casos: pelo painel (\ucr{UC-37}, o caminho preferido) e pelo
\cd{publicar-manifesto.mjs --versao} (\ucr{UC-23}), que continua valendo para
quem opera pelo terminal.
\end{nota}

% ---------------------------------------------------------------------------
\subsection{Visão lógica: classes e objetos}
\label{sec:uml-logica}

\subsubsection{O modelo de domínio da plataforma}

O modelo entidade-relacionamento do Capítulo~\ref{cap:dados} diz como as tabelas
se ligam. O diagrama de classes da Figura~\ref{fig:uml-classes-dominio} diz outra
coisa: que \emph{conceitos} a plataforma tem, e quais regras são restrição do
modelo em vez de detalhe de consulta. Três notações carregam o que o ER não
carrega. O \textbf{atributo derivado} (barra antes do nome) marca os dois
\emph{caches} de \cd{users}: \cd{/escola} é a lotação principal e \cd{/turma} é a
matrícula principal, e quem escreve os dois é \cd{matricularEm}, nunca a tela. A
\textbf{classe abstrata} \cd{Recurso} é a da \adrr{ADR-87}: produto e experiência
são recursos, e o \emph{produto} é o teto das experiências que declara. E as
\textbf{restrições entre chaves} dizem, no próprio elemento, o que é invariável.
Duas delas são o contraste da regra 8 do \arq{AGENTS.md}: ausência de linha é
\emph{bloqueada} em \cd{Liberacao} e \emph{ligada} em \cd{Disponibilidade} ---
salvo o produto em teste interno, que nasce desligado
(\cd{PRODUTOS\_\B DESLIGADOS\_\B POR\_\B PADRAO}, hoje só a \cd{ajuda},
\adrr{ADR-132}). O padrão mora no resolvedor, não numa linha semeada: numa
\emph{stack} nova não há usuário para assinar \cd{definido\_por}.

\begin{figure}[H]
\centering
\fitpage{%
\begin{tikzpicture}[x=1cm, y=1cm]

  % sem hifenizacao: a quebra das restricoes vai a mao
  \hyphenpenalty=10000 \exhyphenpenalty=10000
  % ---- linha 1: escola, turma, periodo
  \node[classe, text width=42mm, anchor=north] (esc) at (2.6,0)
    {\etitulo{Escola}
     \nodepart{two}
     \texttt{nome: String}\\
     \texttt{codigoInep: String [0..1]}\\
     \texttt{tipo: TipoDeEscola}\\
     \texttt{municipio, uf: String [0..1]}
     \nodepart{three}
     \{regional é Escola com\\
     \ \ \texttt{tipo = regional}\}\\
     \{o contrato não é linha:\\
     \ \ é a \emph{stack}\}};

  \node[classe, text width=40mm, anchor=north] (tur) at (8.6,0)
    {\etitulo{Turma}
     \nodepart{two}
     \texttt{nome: String}\\
     \texttt{serie: String}\\
     \texttt{tipo: regular | componente}\\
     \texttt{turno: String [0..1]}
     \nodepart{three}
     \{única por escola, período\\
     \ \ e nome\}\\
     \{fora da RLS: o filtro é a tranca\}};

  \node[classe, text width=34mm, anchor=north] (per) at (14.4,0)
    {\etitulo{PeriodoLetivo}
     \nodepart{two}
     \texttt{tipo, titulo: String}\\
     \texttt{ano: Int}\\
     \texttt{inicio, fim: Date}
     \nodepart{three}
     \{do contrato, sem escola\}};

  % ---- linha 2: lotacao, matricula, liberacao
  \node[classe, text width=42mm, anchor=north] (lot) at (2.6,-5.0)
    {\etitulo{Lotacao}
     \nodepart{two}
     \texttt{papel: estudante | professor}\\
     \texttt{\ \ | gestor\_escola}\\
     \texttt{principal: Boolean}\\
     \texttt{inicio: Date}\\
     \texttt{fim: Date [0..1]}
     \nodepart{three}
     \{única por escola, usuário\\
     \ \ e papel\}};

  \node[classe, text width=40mm, anchor=north] (mat) at (8.6,-5.0)
    {\ctitulo{TabelaDeTenant}{Matricula}
     \nodepart{two}
     \texttt{papel: estudante}\\
     \texttt{\ \ | professor}\\
     \texttt{principal: Boolean}\\
     \texttt{inicio: Date}\\
     \texttt{fim: Date [0..1]}
     \nodepart{three}
     \{no máximo uma principal ativa\\
     \ \ por estudante\}};

  \node[classe, text width=34mm, anchor=north] (lib) at (14.4,-5.0)
    {\ctitulo{TabelaDeTenant}{Liberacao}
     \nodepart{two}
     \texttt{estado: liberada}\\
     \texttt{\ \ | encerrada}\\
     \texttt{liberadaEm: DateTime}\\
     \texttt{encerradaEm [0..1]}
     \nodepart{three}
     \{sem linha = bloqueada\}\\
     \{única por turma e \texttt{expId}\}};

  % ---- linha 3: usuario, disponibilidade, experiencia
  \node[classe, text width=42mm, anchor=north] (usu) at (2.6,-10.1)
    {\etitulo{Usuario}
     \nodepart{two}
     \texttt{username: String}\\
     \texttt{displayName: String}\\
     \texttt{role: Papel}\\
     \texttt{ativo: Boolean}\\
     \texttt{senhaProvisoria: Boolean}\\
     \texttt{/escola: Escola [0..1]}\\
     \texttt{/turma: Turma [0..1]}
     \nodepart{three}
     \{escola obrigatória, salvo\\
     \ \ papéis sem escola\}\\
     \{\texttt{role = gestor\_regional} \texttt{<=>}\\
     \ \ \texttt{regional} não nula\}};

  \node[classe, text width=40mm, anchor=north] (dis) at (8.6,-10.1)
    {\etitulo{Disponibilidade}
     \nodepart{two}
     \texttt{ligado: Boolean}\\
     \texttt{versaoFixa: String [0..1]}\\
     \texttt{motivo: String [0..1]}\\
     \texttt{definidoEm: DateTime}
     \nodepart{three}
     \{sem linha = ligado, salvo\\
     \ \ produto em teste interno\}\\
     \{única por recurso e escopo\}\\
     \{escopo nulo = o contrato\}};

  \node[entidade, text width=34mm, anchor=north] (exp) at (14.4,-10.1)
    {\etitulo{Experiencia}
     \nodepart{two}
     \texttt{expId: String}\\
     \texttt{versaoAtiva: String}\\
     \texttt{series: String [*]}\\
     \texttt{publico: Publico}\\
     \texttt{permissoes: String [*]}};

  % ---- linha 4: papel, recurso, produto
  \node[enumeracao, text width=30mm, anchor=north] (pap) at (2.6,-16.3)
    {\ctitulo{enumeração}{Papel}
     \nodepart{two}
     \texttt{estudante}\\
     \texttt{professor}\\
     \texttt{gestor\_escola}\\
     \texttt{gestor\_regional}\\
     \texttt{gestor\_contrato}\\
     \texttt{admin}};

  \node[entidade, text width=40mm, anchor=north,
        rectangle split part fill={unooInkSoft, unooSurface}] (rec) at (8.6,-15.4)
    {\textcolor{white}{\textbf{\textit{Recurso}}}
     \nodepart{two}
     \texttt{/recursoId: String}\\
     {\color{unooMuted}\texttt{produto:<slug>} ou \texttt{experiencia:<expId>}}};

  \node[entidade, text width=34mm, anchor=north] (pro) at (14.4,-15.4)
    {\etitulo{Produto}
     \nodepart{two}
     \texttt{slug: Produto}\\
     {\color{unooMuted}vocabulário em \texttt{recursos/catalogo.ts}}};

  % ---- associacoes
  % escola -- turma
  \draw[associa] (esc.east |- 0,-1.3) -- (tur.west |- 0,-1.3)
    node[multip, pos=0.1, above]{1} node[multip, pos=0.9, above]{*};
  % turma -- periodo
  \draw[associa] (tur.east |- 0,-1.3) -- (per.west |- 0,-1.3)
    node[multip, pos=0.12, above]{*} node[multip, pos=0.88, above]{1};
  % escola (pai) -- escola (filhas): laco pela direita do cabecalho
  \draw[associa] ([xshift=8mm]esc.north) -- ++(0,0.55) -| ([xshift=4mm]esc.east |- 0,-0.25)
    -- (esc.east |- 0,-0.25);
  \node[papelassoc, anchor=south east] at ([xshift=8mm, yshift=0.55cm]esc.north) {filhas *};
  \node[papelassoc, anchor=south west] at ([xshift=4mm, yshift=0.55cm]esc.north east) {pai 0..1};
  % escola -- lotacao -- usuario
  \draw[associa] (esc.south) -- (lot.north)
    node[multip, pos=0.15, right]{1} node[multip, pos=0.85, right]{*};
  \draw[associa] (lot.south) -- (usu.north)
    node[multip, pos=0.15, right]{*} node[multip, pos=0.85, right]{1};
  % turma -- matricula -- usuario
  \draw[associa] (tur.south) -- (mat.north)
    node[multip, pos=0.15, right]{1} node[multip, pos=0.85, right]{*};
  \draw[associa] (usu.east |- 0,-11.0) -- (5.95,-11.0) |- (mat.west |- 0,-7.0)
    node[multip, pos=0.92, above]{*};
  \node[multip, anchor=south west] at (usu.east |- 0,-11.0) {1};
  % turma -- liberacao -- experiencia
  \draw[associa] (tur.east |- 0,-3.1) -| ([xshift=-9mm]lib.north)
    node[multip, pos=0.06, above]{1} node[multip, pos=0.95, right]{*};
  \draw[associa] (lib.south) -- (exp.north)
    node[multip, pos=0.15, right]{*} node[multip, pos=0.85, right]{1}
    node[papelassoc, midway, left=2pt, align=right]{por \texttt{expId},\\sem FK};
  % usuario -- escola (regional do gestor)
  \draw[associa] (usu.west) -- ++(-0.55,0) |- (esc.west);
  \node[papelassoc, anchor=south east] at (esc.west) {regional 0..1};
  % disponibilidade -- escola (escopo), pelo corredor entre as colunas
  \draw[associa] (dis.west |- 0,-12.0) -- (5.45,-12.0) |- (esc.east |- 0,-3.3);
  \node[papelassoc, anchor=south west] at (esc.east |- 0,-3.3) {escopo 0..1};
  \node[multip, anchor=south east] at (dis.west |- 0,-12.0) {*};
  % disponibilidade -- recurso
  \draw[associa] (dis.south) -- (rec.north)
    node[multip, pos=0.2, right]{*} node[multip, pos=0.8, right]{1};
  % recurso <|-- produto, experiencia (alvo compartilhado)
  \draw[hered] (pro.west |- rec.east) -- (rec.east);
  \draw[associa] (exp.west |- 0,-13.0) -| (11.55,0 |- rec.east);
  % produto -- experiencia (o teto)
  \draw[associa] (pro.north) -- (exp.south)
    node[multip, pos=0.15, right]{0..1} node[multip, pos=0.85, right]{*}
    node[papelassoc, midway, left=2pt]{teto};
  % usuario -- papel
  \draw[depend] (usu.south) -- (pap.north);

\end{tikzpicture}}
\caption[Classes: o modelo de domínio da plataforma]{O modelo de domínio da
  plataforma. \estereo{TabelaDeTenant} marca só as duas tabelas com RLS forçada;
  \cd{users}, \cd{turmas} e \cd{lotacoes} têm \cd{escola\_id} mas ficam fora da
  RLS, porque o \emph{login} as lê antes de saber a escola --- nelas o filtro da
  aplicação é a única tranca (Capítulo~\ref{cap:tenancy}). Fontes:
  \arq{apps/api/prisma/schema.prisma}, os CHECKs das migrações
  \arq{modelo\_de\_papeis}, \arq{nivel\_regional} e \arq{liberacoes\_por\_turma},
  \arq{rostering/matricula.ts}, \arq{recursos/regras.ts} e
  \arq{recursos/catalogo.ts}.}
\label{fig:uml-classes-dominio}
\end{figure}

\subsubsection{A porta da IA: estratégia, decorador e quota}

A regra 5 do \arq{AGENTS.md} --- IA só por \arq{src/ai}, quota antes e depois ---
tem uma forma de projeto, e é ela que a Figura~\ref{fig:uml-classes-ia} mostra.
São dois padrões conhecidos. A \textbf{estratégia} é a interface \cd{AiProvider}
com duas implementações trocáveis por \emph{env} (\env{AI\_PROVIDER}: \cd{vertex}
ou \cd{mock}). O \textbf{decorador} é o \cd{MedidorDeProvider}, que implementa a
mesma interface e embrulha a implementação real para medir latência, \emph{tokens}
e \cd{finishReason} (Capítulo~\ref{cap:observabilidade}). Quem monta os dois é a
fábrica do \cd{AiModule}, e nenhum consumidor sabe qual veio: todos injetam o
\emph{token} de injeção \cd{AI\_PROVIDER}. Não existe \cd{AiService} de fachada
--- a busca no código não o encontra, e o diagrama não o inventa. Cada
consumidor chama o \emph{provider} e o \cd{QuotaService} por conta própria, e é por
isso que a quota é regra de revisão e não de tipo: nada impede um \emph{service}
novo de injetar o \emph{provider} e esquecer o \cd{temSaldo}.

Desde 2026-10-02 a interface tem uma sexta operação, \cd{conversarEmFluxo}: a
única geração em \emph{stream} do \emph{provider} (\cd{streamGenerateContent?alt=sse}
no Vertex), sem ferramentas e sem \cd{validar}, porque num \emph{stream} o texto
já saiu quando a validação poderia recusá-lo. Ela existe para a Ajuda com avatar
(\adrr{ADR-132}), e a Ajuda é também o único consumidor que \textbf{não} chama o
\cd{QuotaService}, por decisão: o custo dela é limitado em minutos de vídeo
(\cd{ajuda\_\B limites}), e os \emph{tokens} continuam medidos pelo
\cd{MedidorDeProvider}. O vídeo em si fica fora desta figura, numa terceira
porta de \arq{src/ai}, a \cd{AvatarAoVivoPort} (Figura~\ref{fig:uml-seq-ajuda}).

A mesma fábrica monta uma segunda porta, à direita da figura: a do agente de
IA do professor (\cd{AgenteProfessorPort}, \emph{token}
\cd{AGENTE\_PROFESSOR}; contrato em \arq{docs/regras/agente-professor.md}).
Ela é porta própria, e não método do \cd{AiProvider}, porque o modelo, as
ferramentas e a instrução moram no agente (Python/ADK no Vertex AI Agent
Engine) e a API só conduz a sessão. O \cd{ClienteAgentEngine} é a única peça
da API que conhece o Agent Engine; o mesmo arquivo traz o
\cd{AgenteProfessorMock} da suíte e2e e o \cd{AgenteProfessorDesligado}, que é
o padrão em qualquer ambiente. O \cd{MedidorDeProvider} não o embrulha, mas a
quota vale do mesmo jeito: o \cd{AgenteProfessorService} chama \cd{temSaldo}
antes do turno e \cd{consumir} depois, com o perfil \cd{professor}.

Desde a \adrr{ADR-161} o decorador tem uma segunda saída, que a figura não
desenha: além das métricas, ele entrega cada chamada ao \cd{RegistroDeConsumo}
(\arq{ai/consumo-de-ia.ts}), que soma o que foi pago em
\cd{ia\_consumo\_diario} para o painel de custos. É o mesmo ponto de
captura, e por isso a quota e o custo nunca divergem sobre o que é uma
chamada: o \cd{consumir} da quota (DEPOIS) e a linha de consumo saem da mesma
contagem do \emph{vendor} --- no \cd{generateJson}, do mesmo
\cd{aoUsarTokens}, que agora leva também a divisão entre entrada e saída. A
quota continua sendo o que barra; o registro só mede, e nunca derruba a
chamada (Seção~\ref{sec:ia-custos}).

\begin{figure}[H]
\centering
\fitpage{%
\begin{tikzpicture}[x=1cm, y=1cm]
  % sem hifenizacao: as caixas sao estreitas, a quebra vai a mao
  \hyphenpenalty=10000 \exhyphenpenalty=10000
  % quatro colunas (A, B, C, D) e caixas que ocupam duas (AB, CD)
  \def\xA{2.0}\def\xB{6.25}\def\xC{10.5}\def\xD{14.75}
  \def\xAB{4.125}\def\xBC{8.375}\def\xCD{12.625}
  \def\cm#1{{\color{unooMuted}#1}}

  % ---- linha 1: a quota e as duas tabelas que ela le e escreve
  \node[tabela, text width=27mm, anchor=north] (tq) at (1.55,0)
    {\etitulo{ia\_quotas}
     \nodepart{two}
     \pk{perfil} \texttt{PK}\\
     \texttt{tokens\_por\_dia}\\
     \cm{sem escola: do contrato}};
  \node[classe, text width=72mm, anchor=north] (quota) at (\xBC,0)
    {\etitulo{QuotaService}
     \nodepart{two}
     \texttt{- prisma: PrismaService}
     \nodepart{three}
     \texttt{+ limite(perfil): Promise<number>}\\
     \texttt{+ usadoHoje(escolaId, usuarioId)}\\
     \texttt{\ \ \ \ : Promise<number>}\\
     \texttt{+ temSaldo(escolaId, usuarioId,}\\
     \texttt{\ \ \ \ perfil = "estudante"): Promise<boolean>}\\
     \texttt{+ consumir(escolaId, usuarioId, tokens)}\\
     \texttt{\ \ \ \ : Promise<number>}};
  \node[tabela, text width=27mm, anchor=north] (tu) at (15.2,0)
    {\ctitulo{TabelaDeTenant}{ia\_uso\_diario}
     \nodepart{two}
     \pk{escola\_id,}\\
     \pk{usuario\_id, dia}\\
     \texttt{tokens\_usados}};

  % ---- linha 2: os consumidores
  \node[classe, text width=34mm, anchor=north] (ctrl) at (\xA,-4.5)
    {\ctitulo{controller}{AiController}
     \nodepart{two}
     \texttt{- provider: AiProvider}\\
     \texttt{- quota: QuotaService}
     \nodepart{three}
     \texttt{+ gerar(corpo): \{texto\}}\\
     \cm{\texttt{POST /ai/gerar},}\\
     \cm{só \texttt{OidcGuard}}};
  \node[classe, text width=34mm, anchor=north] (dom) at (\xB,-4.5)
    {\ctitulo{service}{PlanosService}
     \nodepart{two}
     \cm{e mais treze módulos}\\
     \cm{de domínio}
     \nodepart{three}
     \texttt{+ gerar(\ldots)}\\
     \cm{\texttt{temSaldo} \seta{} \emph{provider}}\\
     \cm{\seta{} \texttt{consumir}}};
  \node[classe, text width=34mm, anchor=north] (ags) at (\xD,-4.5)
    {\ctitulo{service}{AgenteProfessorService}
     \nodepart{two}
     \texttt{- agente:}\\
     \texttt{\ \ AgenteProfessorPort}\\
     \texttt{- quota: QuotaService}
     \nodepart{three}
     \texttt{+ temSaldo(quem)}\\
     \texttt{+ transmitir(\ldots)} \cm{SSE}};
  \node[actnote, text width=32mm, anchor=north] (nota) at (\xC,-4.5)
    {\textbf{O que o tipo não garante:} nada obriga a chamar
     \texttt{temSaldo}, e entre ele e o \texttt{consumir} não há trava:
     duas chamadas simultâneas podem passar o teto em uma resposta
     cada.};

  % ---- linha 3: as duas portas e o tipo de resultado
  \node[interface, text width=72mm, anchor=north] (prov) at (\xAB,-8.7)
    {\ctitulo{interface}{AiProvider}
     \nodepart{two}
     \cm{injetado pelo \emph{token} \texttt{AI\_PROVIDER}}
     \nodepart{three}
     \texttt{+ generate(req): Promise<AiResult>}\\
     \texttt{+ generateJson<T>(req): Promise<T>}\\
     \texttt{+ generateImage(prompt, opcoes?)}\\
     \texttt{\ \ \ \ : Promise<Buffer>}\\
     \texttt{+ generateSpeech(req): Promise<Buffer>}\\
     \texttt{+ conversar(req): Promise<AiResult>}\\
     \texttt{+ conversarEmFluxo(req): AsyncIterable}};
  \node[entidade, text width=32mm, anchor=north,
        rectangle split part fill={unooMuted, unooSurface}] (tipos) at (10.75,-8.7)
    {\ctitulo{tipo de dado}{AiResult}
     \nodepart{two}
     \texttt{texto: string}\\
     \texttt{tokensUsados: number}\\
     \texttt{tokensEntrada?,}\\
     \texttt{tokensSaida?}\\
     \texttt{finishReason?}\\
     \texttt{voltas?, tentativas?}};
  \node[interface, text width=34mm, anchor=north] (port) at (\xD,-8.7)
    {\ctitulo{interface}{AgenteProfessorPort}
     \nodepart{two}
     \texttt{+ modo: ModoDoAgente}\\
     \cm{injetada pelo \emph{token}}\\
     \cm{\texttt{AGENTE\_PROFESSOR}}
     \nodepart{three}
     \texttt{+ criarSessao(\ldots)}\\
     \texttt{+ lerSessao(\ldots)}\\
     \texttt{+ apagarSessao(\ldots)}\\
     \texttt{+ renovarSessao?(\ldots)}\\
     \texttt{+ enviar(\ldots)}\\
     \texttt{\ \ : AsyncIterable<\ldots>}};
  \path (port.south) ++(0,-0.55) coordinate (r4);

  % ---- linha 4: as implementacoes
  \node[classe, text width=34mm, anchor=north] (vert) at (\xA,0 |- r4)
    {\etitulo{VertexProvider}
     \nodepart{two}
     \texttt{- auth: GoogleAuth}\\
     \cm{ADC, sem chave em \emph{env}}
     \nodepart{three}
     \cm{as seis operações por}\\
     \cm{\texttt{:generateContent};}\\
     \cm{o fluxo, por}\\
     \cm{\texttt{:streamGenerateContent}}};
  \node[classe, text width=34mm, anchor=north] (mock) at (\xB,0 |- r4)
    {\etitulo{MockProvider}
     \nodepart{two}
     \cm{determinístico}
     \nodepart{three}
     \texttt{generate}, \texttt{conversar},\\
     \texttt{conversarEmFluxo}:\\
     respondem; os outros três:\\
     \textbf{lançam}};
  \node[classe, text width=34mm, anchor=north] (med) at (\xC,0 |- r4)
    {\ctitulo{decorador}{MedidorDeProvider}
     \nodepart{two}
     \texttt{- real: AiProvider}\\
     \texttt{- metricas: Metricas}\\
     \texttt{- modelos:}\\
     \texttt{\ \ ModelosDoMedidor}
     \nodepart{three}
     \cm{as seis operações, cada}\\
     \cm{uma por \texttt{medir(op,}}\\
     \cm{\texttt{modelo, fn)}}};
  \node[classe, text width=34mm, anchor=north] (cli) at (\xD,0 |- r4)
    {\etitulo{ClienteAgentEngine}
     \nodepart{two}
     \cm{modo \texttt{agent-engine}}\\
     \cm{(ADC) ou \texttt{local} (\emph{shim})}
     \nodepart{three}
     \cm{\texttt{:query} e}\\
     \cm{\texttt{:streamQuery?alt=sse}}\\
     \cm{no Agent Engine}};
  \path (med.south) ++(0,-0.9) coordinate (r5);

  % ---- linha 5: o erro, a fabrica e a nota da fabrica
  \node[entidade, text width=30mm, anchor=north] (erro) at (\xB,0 |- r5)
    {\etitulo{ErroDeGeracao}
     \nodepart{two}
     \texttt{extends Error}};
  \node[classe, text width=72mm, anchor=north] (mod) at (\xCD,0 |- r5)
    {\ctitulo{módulo @Global}{AiModule}
     \nodepart{two}
     \texttt{exports: [AI\_PROVIDER, QuotaService,}\\
     \texttt{\ \ RegistroDeConsumo, AGENTE\_PROFESSOR,}\\
     \texttt{\ \ AVATAR\_AO\_VIVO, VOZ\_AO\_VIVO]}
     \nodepart{three}
     \texttt{+ useFactory(config, metricas, registro)}\\
     \texttt{\ \ \ \ : AiProvider} \cm{(\texttt{AI\_PROVIDER})}\\
     \texttt{+ useFactory(config): AgenteProfessorPort}\\
     \texttt{\ \ \ \ }\cm{\texttt{= criarAgenteProfessor(\{modo, \ldots\})}}};
  \node[actnote, text width=39mm, anchor=north] (fab) at (2.2,0 |- r5)
    {\textbf{A fábrica, no código:}\\
     \texttt{new MedidorDeProvider(}\\
     \texttt{\ env === "mock"}\\
     \texttt{\ ? new MockProvider()}\\
     \texttt{\ : new VertexProvider(),}\\
     \texttt{\ metricas, modelos, \ldots)}.\\
     Trocar de fornecedor é um arquivo novo em \texttt{src/ai} e uma
     linha aqui.};

  % ---- quem usa a quota (pela faixa entre as linhas 1 e 2)
  \def\yg{-4.05}
  \draw[depend] (ctrl.north) |- (4.9,\yg) -- (4.9,0 |- quota.south);
  \draw[depend] (dom.north) -- (dom.north |- quota.south);
  \draw[depend] (ags.north) |- (11.9,\yg) -- (11.9,0 |- quota.south);
  \draw[depend] (quota.west |- tq.east) -- (tq.east)
    node[cxlabel, midway, above=1pt]{lê};
  \draw[depend] (quota.east |- tu.west) -- (tu.west)
    node[cxlabel, midway, above=1pt]{\emph{upsert}};

  % ---- quem usa cada porta
  \draw[depend] (ctrl.south) -- (ctrl.south |- prov.north);
  \draw[depend] (dom.south) -- (dom.south |- prov.north)
    node[cxlabel, midway, right=2pt]{\texttt{@Inject(AI\_PROVIDER)}};
  \draw[depend] (ags.south) -- (port.north)
    node[cxlabel, midway, left=2pt]{\texttt{@Inject(AGENTE\_PROFESSOR)}};
  \draw[depend] (prov.east |- 0,-9.0) -- (tipos.west |- 0,-9.0)
    node[cxlabel, midway, above=1pt, inner sep=0.8pt]{\estereo{usa}};

  % ---- realizacoes, e o decorador que embrulha o real
  \draw[realiza] (vert.north) -- (vert.north |- prov.south);
  \draw[realiza] ([xshift=-5mm]mock.north) -- ([xshift=-5mm]mock.north |- prov.south);
  \path (prov.south) ++(0,-0.35) coordinate (yB) ++(0,-0.4) coordinate (yA);
  \draw[agrega] (9.3,0 |- med.north) |- (6.95,0 |- yA) -- (6.95,0 |- prov.south)
    node[papelassoc, pos=0.5, left=1pt]{\texttt{real} 1};
  \draw[realiza] (10.1,0 |- med.north) |- (7.65,0 |- yB) -- (7.65,0 |- prov.south);
  \draw[realiza] (cli.north) -- (port.south);
  \draw[depend] (mock.south) -- (erro.north)
    node[cxlabel, midway, right=2pt]{\estereo{lança}};

  % ---- a fabrica cria o decorador e o cliente do agente
  \draw[depend] (\xC,0 |- mod.north) -- (med.south)
    node[cxlabel, midway, right=2pt]{\estereo{cria}};
  \draw[depend] (\xD,0 |- mod.north) -- (cli.south)
    node[cxlabel, midway, right=2pt]{\estereo{cria}};

\end{tikzpicture}}
\caption[Classes: a porta da IA]{A porta da IA como diagrama de classes de
  projeto: estratégia (\cd{AiProvider} com duas implementações) e decorador
  (\cd{MedidorDeProvider}), montados pela fábrica do \cd{AiModule} e
  injetados pelo \emph{token} \cd{AI\_PROVIDER}. A quota é um serviço à parte,
  que cada consumidor chama --- menos a Ajuda com avatar, limitada por minutos
  (\adrr{ADR-132}). À direita, a porta do agente de IA do professor, que a
  mesma fábrica monta. Fontes: \arq{ai/provider.ts}, \arq{ai/ai.module.ts},
  \arq{ai/medidor.provider.ts}, \arq{ai/quota.service.ts},
  \arq{ai/ai.controller.ts}, \arq{ai/agente-professor.port.ts},
  \arq{ai/agent-engine.client.ts}.}
\label{fig:uml-classes-ia}
\end{figure}

\subsubsection{Um instantâneo da disponibilidade}

As três regras da disponibilidade --- o escopo mais específico ganha, o produto é
teto, ausência de linha é ligado (salvo produto em teste interno) --- são fáceis de enunciar e fáceis de errar
quando se combinam. O diagrama de objetos da Figura~\ref{fig:uml-objetos} congela
uma rede com uma regional, duas escolas e quatro linhas em \cd{disponibilidades},
e a Tabela~\ref{tab:uml-objetos} diz o que \cd{exigirRecursoDisponivel} decide
em cada combinação. Os valores são de exemplo; os campos, os códigos de
experiência e os \emph{slugs} de produto são os reais.

\begin{figure}[H]
\centering
\fitwidth{%
\begin{tikzpicture}[x=1cm, y=1cm]
  \tikzset{ob/.style={objeto, text width=38mm}}

  % ---- a arvore da rede
  \node[ob, anchor=north] (norte) at (8.45,0)
    {\underline{\textbf{norte : Escola}}
     \nodepart{two}
     \texttt{tipo = "regional"}\\
     \texttt{nome = "Regional Norte"}};
  \node[ob, anchor=north] (ef) at (6.2,-2.9)
    {\underline{\textbf{ef : Escola}}
     \nodepart{two}
     \texttt{tipo = "escola"}\\
     \texttt{nome = "EE Cora Coralina"}};
  \node[ob, anchor=north] (ee) at (10.7,-2.9)
    {\underline{\textbf{ee : Escola}}
     \nodepart{two}
     \texttt{tipo = "escola"}\\
     \texttt{nome = "EE Monteiro Lobato"}};

  % ---- as quatro linhas de disponibilidade
  \node[ob, anchor=north] (d3) at (1.7,-2.9)
    {\underline{\textbf{d3 : Disponibilidade}}
     \nodepart{two}
     \texttt{recursoId =}\\
     \texttt{"experiencia:exp06ef2b"}\\
     \texttt{ligado = false}\\
     \texttt{motivo = "adiar"}};
  \node[ob, anchor=north] (d4) at (15.2,-2.9)
    {\underline{\textbf{d4 : Disponibilidade}}
     \nodepart{two}
     \texttt{recursoId =}\\
     \texttt{"produto:plano-de-aula"}\\
     \texttt{ligado = false}\\
     \texttt{motivo = "não contratado"}};
  \node[ob, anchor=north] (d1) at (1.7,-6.9)
    {\underline{\textbf{d1 : Disponibilidade}}
     \nodepart{two}
     \texttt{recursoId =}\\
     \texttt{"experiencia:exp06ef2b"}\\
     \texttt{escola = nulo} {\color{unooMuted}(contrato)}\\
     \texttt{ligado = false}};
  \node[ob, anchor=north] (d2) at (6.2,-6.9)
    {\underline{\textbf{d2 : Disponibilidade}}
     \nodepart{two}
     \texttt{recursoId =}\\
     \texttt{"experiencia:exp06ef2b"}\\
     \texttt{ligado = true}\\
     \texttt{motivo = "piloto"}};

  % ---- sala, recursos
  \node[ob, anchor=north] (t6a) at (10.7,-6.9)
    {\underline{\textbf{t6a : Turma}}
     \nodepart{two}
     \texttt{nome = "6A"}\\
     \texttt{serie = "6EF"}};
  \node[ob, anchor=north] (paula) at (15.2,-6.9)
    {\underline{\textbf{aula : Produto}}
     \nodepart{two}
     \texttt{slug = "plano-de-aula"}};
  \node[ob, anchor=north] (e1) at (3.95,-11.0)
    {\underline{\textbf{e1 : Experiencia}}
     \nodepart{two}
     \texttt{expId = "exp06ef2b"}\\
     \texttt{produto = "entre-redes"}\\
     \texttt{series = ["6EF"]}};
  \node[ob, anchor=north] (l1) at (10.7,-11.0)
    {\underline{\textbf{l1 : Liberacao}}
     \nodepart{two}
     \texttt{estado = "liberada"}};
  \node[ob, anchor=north] (e2) at (15.2,-11.0)
    {\underline{\textbf{e2 : Experiencia}}
     \nodepart{two}
     \texttt{expId =}\\
     \texttt{"app-90-plano-de-aula"}\\
     \texttt{publico = "professor"}};

  % ---- links (instancias das associacoes da Figura de classes)
  \draw[objetolink] (ef.north) -- (norte.south west) node[papelassoc, pos=0.45, right=2pt]{pai};
  \draw[objetolink] (ee.north) -- (norte.south east) node[papelassoc, pos=0.45, left=2pt]{pai};
  \draw[objetolink] (d3.north) |- (5.2,-2.4) -- (5.2,0 |- ef.north);
  \node[papelassoc, anchor=south] at (3.45,-2.4) {escopo};
  \draw[objetolink] (d4.north) |- (11.7,-2.4) -- (11.7,0 |- ee.north);
  \node[papelassoc, anchor=south] at (13.45,-2.4) {escopo};
  \draw[objetolink] (d2.north) -- ++(0,0.35) -| (norte.south)
    node[papelassoc, pos=0.62, left=1pt]{escopo};
  \draw[objetolink] (t6a.north) -- (ee.south) node[papelassoc, midway, right=2pt]{escola};
  \draw[objetolink] (l1.north) -- (t6a.south) node[papelassoc, midway, right=2pt]{turma};
  \draw[objetolink] (l1.west) -- (e1.east |- l1.west) node[papelassoc, midway, above]{\texttt{expId}};
  \draw[objetolink] (2.8,0 |- d1.south) -- (2.8,0 |- e1.north) node[papelassoc, midway, right=2pt]{recurso};
  \draw[objetolink] (5.1,0 |- d2.south) -- (5.1,0 |- e1.north) node[papelassoc, midway, right=2pt]{recurso};
  \draw[objetolink] (d3.west) -- ++(-0.35,0) |- (e1.west) node[papelassoc, pos=0.9, above]{recurso};
  \draw[objetolink] (d4.south) -- (paula.north) node[papelassoc, midway, right=2pt]{recurso};
  \draw[objetolink] (e2.north) -- (paula.south) node[papelassoc, midway, right=2pt]{produto (teto)};

\end{tikzpicture}}
\caption[Objetos: um instantâneo da disponibilidade]{Diagrama de objetos: uma rede
  com uma regional, duas escolas e quatro linhas de disponibilidade. Cada caixa é
  uma instância (nome sublinhado) e cada linha é um \emph{link}, instância de uma
  associação da Figura~\ref{fig:uml-classes-dominio}. \texttt{d1} não tem
  \emph{link} de escopo: escola nula é o contrato inteiro.}
\label{fig:uml-objetos}
\end{figure}

\begin{tabularx}{\linewidth}{@{}L{42mm}L{20mm}X@{}}
\caption{O que o resolvedor decide no instantâneo da Figura~\ref{fig:uml-objetos}}
\label{tab:uml-objetos}\\
\toprule
\headrow \thd{Recurso, na escola} & \thd{Decisão} & \thd{Por quê (\texttt{recursos/regras.ts})}\\
\midrule
\endfirsthead
\toprule
\headrow \thd{Recurso, na escola} & \thd{Decisão} & \thd{Por quê}\\
\midrule
\endhead
\cd{exp06ef2b} em \texttt{ee} & ligada & \texttt{ee} não tem linha; a regional
  tem \texttt{d2}, que vence a do contrato (\texttt{d1}).\\
\zebra \cd{exp06ef2b} em \texttt{ef} & desligada & \texttt{d3} é da própria escola
  e vence a regional: o mais específico ganha, nos dois sentidos.\\
\cd{exp06ef2b} numa escola de outra regional & desligada & Nem escola nem
  regional têm linha; vale \texttt{d1}, do contrato. Um estudante com a missão
  liberada recebe 403 \cd{recurso-indisponivel} --- a liberação \texttt{l1} não
  muda nada, porque disponibilidade é pré-condição.\\
\zebra \cd{app-90-plano-de-aula} em \texttt{ee} & desligada & \texttt{d4} desliga o
  \emph{produto} na escola, e o produto é teto: nenhuma linha de experiência
  religaria. O professor lê \enquote{Este recurso não está disponível para a sua
  rede.}\\
\cd{app-90-plano-de-aula} em \texttt{ef} & ligada & Nenhuma linha em nenhum
  escopo, nem do produto nem da experiência.\\
\zebra \cd{exp05ef1b} em qualquer escola & ligada & Ausência de linha é
  \emph{ligado}: o interruptor é mecanismo de exceção.\\
\cd{produto:ajuda} em qualquer escola & desligada & A exceção da exceção: sem
  linha vale o padrão do produto, e a \cd{ajuda} está em
  \cd{PRODUTOS\_\B DESLIGADOS\_\B POR\_\B PADRAO} (teste interno,
  \adrr{ADR-132}). Uma linha \cd{ligado = true} na escola, na regional ou no
  contrato a liga ali; tirar o \emph{slug} da lista muda o padrão de todas as
  redes de uma vez.\\
\bottomrule
\end{tabularx}

% ---------------------------------------------------------------------------
\subsection{Visão de desenvolvimento: pacotes, componentes e estrutura composta}
\label{sec:uml-desenvolvimento}

\subsubsection{Quem pode importar quem}

O diagrama de pacotes da Figura~\ref{fig:uml-pacotes} tem uma pergunta só: que
dependência de código é permitida. As setas saem das declarações reais --- os
\arq{package.json} do \emph{workspace} e os \cd{import} da API --- e duas
ausências valem tanto quanto as setas. \textbf{A API não importa pacote nenhum do
\emph{workspace}}: ela não depende de \cd{@edu/sdk}, e o que a SDK envia (o
lote de telemetria, por exemplo) é validado de novo do lado da API, com o Zod
dela. E
\textbf{hub e API não se importam}: falam por HTTP, pelo BFF, e um tipo
compartilhado entre os dois seria o primeiro passo para o hub assumir regra de
domínio.

\begin{figure}[H]
\centering
\fitwidth{%
\begin{tikzpicture}[x=1cm, y=1cm]
  \tikzset{pk/.style={font=\scriptsize, align=left, anchor=north west, inner sep=2pt}}
  \hyphenpenalty=10000 \exhyphenpenalty=10000

  % ---- o monorepo; os pacotes levam o nome qualificado (apps/, packages/)
  \pacote[fill=unooPaper]{mono}{(0,0)}{(12.4,-18.75)}{\estereo{monorepo} plataforma-unooria}

  % ---- linha 1: a api e as quatro faixas
  \pacote{api}{(0.4,-0.95)}{(12.0,-8.4)}{apps/api}
  \node[pk] at (0.45,-1.0) {\texttt{@edu/api} --- NestJS + Prisma};
  \pacote[fill=unooSurface]{ltr}{(0.8,-2.0)}{(9.0,-2.75)}{transversal}
  \node[pk, text width=80mm] at (0.85,-2.05) {\texttt{auth/} \ \texttt{tenancy/} \ \texttt{observabilidade/} \ \texttt{config/}};
  \pacote[fill=unooSurface]{lpl}{(0.8,-3.75)}{(9.0,-4.95)}{plataforma}
  \node[pk, text width=80mm] at (0.85,-3.8) {\texttt{usuarios/ rostering/ salas/ equipes/ evidencias/}\\\texttt{recursos/ experiencias/ telemetry/ analise/ comum/} \ldots};
  \pacote[fill=unooSurface]{lin}{(0.8,-5.75)}{(9.0,-6.5)}{infraestrutura}
  \node[pk, text width=80mm] at (0.85,-5.8) {\texttt{prisma/} \ \texttt{ai/} \ \texttt{object-storage/} \ \texttt{assets/}};
  \pacote[fill=unooGreen!10]{ldo}{(0.8,-7.3)}{(9.0,-8.0)}{domínio \{nenhum módulo importa outro\}}
  \node[pk, text width=80mm] at (0.85,-7.35) {\texttt{motor/ geracao/ vila/ unoogpt/ quem-mandou-isso/} \ldots};

  % dependencias entre faixas, na pista da direita (x = 9.4 a 10.4)
  \draw[pacotedep] (ldo.east |- 0,-7.85) -- (10.4,-7.85) -- (10.4,-4.1) -- (lpl.east |- 0,-4.1);
  \draw[pacotedep] (ldo.east |- 0,-7.55) -- (9.9,-7.55) -- (9.9,-6.3) -- (lin.east |- 0,-6.3);
  \draw[pacotedep] (lpl.east |- 0,-4.6) -- (9.4,-4.6) -- (9.4,-5.95) -- (lin.east |- 0,-5.95);
  \node[cxlabel, rotate=90, anchor=north] at (10.4,-6.0) {\estereo{import}};
  \draw[pacotedep] (lpl.north -| 4.9,0) -- (ltr.south -| 4.9,0)
    node[cxlabel, midway, right=2pt]{\estereo{import} --- toda faixa usa a transversal};

  % ---- linha 2: o hub e o agente do professor
  \pacote{hub}{(0.4,-9.35)}{(5.0,-11.6)}{apps/hub}
  \node[pk, text width=42mm] at (0.45,-9.4)
    {\texttt{@edu/hub}\\Next.js, App Router\\rotas por papel, BFF, JWKS,\\montagem do \emph{embed}};
  \pacote{agp}{(6.2,-9.35)}{(12.0,-11.15)}{apps/agente-professor}
  \node[pk, text width=56mm] at (6.25,-9.4)
    {Python + Google ADK, no\\Vertex AI Agent Engine\\{\color{unooMuted}não importa nem é importado:}\\{\color{unooMuted}só a API o chama, por HTTP}};

  % ---- linha 3: os pacotes do workspace e o externo
  \pacote{sdk}{(0.4,-12.8)}{(5.0,-14.9)}{packages/sdk}
  \node[pk, text width=44mm] at (0.45,-12.85)
    {\texttt{@edu/sdk}\\TypeScript + Zod\\ponte \texttt{postMessage} v2, cliente\\HTTP, telemetria em lote};
  \pacote{ui}{(6.2,-12.8)}{(10.8,-14.9)}{packages/ui}
  \node[pk, text width=44mm] at (6.25,-12.85)
    {\texttt{@edu/ui}\\as oito pontes do DS, uma\\por seção: \texttt{'use client'},\\nomes sem colisão};
  \pacote[fill=unooPaper, draw=unooInkSoft]{ds}{(13.2,-12.8)}{(17.2,-14.9)}{\estereo{externo} npm}
  \node[pk, text width=38mm] at (13.25,-12.85)
    {\texttt{@unooria/design-system}\\repo \texttt{storybook-unooria}\\componentes + CSS\\de seção};

  % ---- linha 4: as experiencias
  \pacote{exps}{(0.4,-16.1)}{(7.4,-18.35)}{apps/experiencias}
  \node[pk, text width=68mm] at (0.45,-16.15)
    {\texttt{\_template} \ \texttt{exp01ef1b} \ldots{} \texttt{exp08ef4b}\\\texttt{app-70} \ \texttt{app-90} \ \texttt{app-91} \ \texttt{app-92}\\{\color{unooMuted}13 \emph{bundles} Vite publicados;}\\{\color{unooMuted}\texttt{\_template} e \texttt{backup-*} ficam fora do \emph{deploy}}};

  % ---- o hub chega aos pacotes por cima, as experiencias por baixo
  \draw[pacotedep] (hub.south -| 2.2,0) -- (sdk.north -| 2.2,0)
    node[cxlabel, midway, right=2pt]{\estereo{import}};
  \draw[pacotedep] (hub.south -| 4.2,0) |- (7.6,-12.2) -- (ui.north -| 7.6,0);
  \node[cxlabel, anchor=south] at (5.9,-12.2) {\estereo{import}};
  \draw[pacotedep] (hub.east |- 0,-11.45) -- (15.2,-11.45) -- (ds.north -| 15.2,0);
  \node[cxlabel, anchor=north] at (9.1,-11.45) {\estereo{use} só CSS da seção e fontes};
  \draw[pacotedep] (exps.north -| 2.2,0) -- (sdk.south -| 2.2,0)
    node[cxlabel, midway, right=2pt]{\estereo{import}};
  \draw[pacotedep] (exps.north -| 6.8,0) -- (ui.south -| 6.8,0)
    node[cxlabel, midway, right=2pt]{\estereo{import}};
  \draw[pacotedep] (exps.east |- 0,-17.2) -- (15.2,-17.2) -- (ds.south -| 15.2,0);
  \node[cxlabel, anchor=south] at (10.1,-17.2) {\estereo{use} só CSS da seção e fontes};
  \draw[pacotedep] (ui.east |- 0,-13.85) -- (ds.west |- 0,-13.85)
    node[cxlabel, midway, above=1pt]{\estereo{import}};

\end{tikzpicture}}
\caption[Pacotes: quem pode importar quem]{Diagrama de pacotes do
  \emph{monorepo}. As setas tracejadas são dependências de código declaradas: o
  hub e as experiências importam a SDK e os componentes pela ponte do DS, e o CSS
  da seção direto do pacote publicado; a API não importa nada do
  \emph{workspace}. A \texttt{app-70} importa só a SDK: é jogo em \emph{canvas},
  sem DOM de interface. O \texttt{agente-professor} (Python) fica fora do grafo
  de \texttt{import}: só a API o chama, por HTTP. No domínio, duas experiências
  conversam por tabela de contrato ou por \emph{endpoint} de fronteira, nunca
  por \texttt{import}. Fontes: os \arq{package.json} do \emph{workspace} e a faixa de cada
  módulo na Figura~\ref{fig:c4l3api}.}
\label{fig:uml-pacotes}
\end{figure}

\subsubsection{Quem fornece e quem requer}

O C4 nível 2 (Figura~\ref{fig:c4l2}) desenha as mesmas peças e diz quem chama
quem. O diagrama de componentes da Figura~\ref{fig:uml-componentes} diz a
\emph{direção da dependência de interface}: o pirulito é a interface que o
componente \textbf{fornece}, o soquete é a que ele \textbf{requer}, e o encaixe dos
dois é um conector de montagem. Duas direções surpreendem quem conhece só o
C4. \textbf{O JWKS é fornecido pelo hub e requerido pela API}: a API depende do
hub para validar qualquer \emph{token}, e é por isso que trocar para o WSO2 é
trocar \env{OIDC\_JWKS\_URI} e \env{OIDC\_ISSUER} sem tocar o \emph{guard}. E
\textbf{a ingestão da telemetria é fornecida pela API e requerida pelo Pub/Sub}:
não há \emph{worker}; é a assinatura \emph{push} que chama
\rota{POST /telemetry/ingest}, autenticada por ID \emph{token} do Google.

\begin{figure}[H]
\centering
\fitwidth{%
\begin{tikzpicture}[x=1cm, y=1cm]
  \hyphenpenalty=10000 \exhyphenpenalty=10000
  % conectores de montagem: bola do lado de quem fornece, soquete do de quem
  % requer; o soquete abre para o lado do fornecedor
  \def\montBaixo#1#2#3#4{% x, borda de quem requer (acima), borda do fornecedor (abaixo), y da bola
    \draw[associa] (#1,#3) -- (#1,#4-0.13);
    \draw[associa] (#1,#2) -- (#1,#4+0.22);
    \bola{(#1,#4)} \soquete{(#1,#4)}{-90}}
  \def\montCima#1#2#3#4{% x, borda do fornecedor (acima), borda de quem requer (abaixo), y da bola
    \draw[associa] (#1,#2) -- (#1,#4+0.13);
    \draw[associa] (#1,#3) -- (#1,#4-0.22);
    \bola{(#1,#4)} \soquete{(#1,#4)}{90}}
  \def\montE#1#2#3#4{% y, borda do fornecedor (esquerda), borda de quem requer (direita), x da bola
    \draw[associa] (#2,#1) -- (#4-0.13,#1);
    \draw[associa] (#3,#1) -- (#4+0.22,#1);
    \bola{(#4,#1)} \soquete{(#4,#1)}{180}}
  \def\montD#1#2#3#4{% y, borda de quem requer (esquerda), borda do fornecedor (direita), x da bola
    \draw[associa] (#3,#1) -- (#4+0.13,#1);
    \draw[associa] (#2,#1) -- (#4-0.22,#1);
    \bola{(#4,#1)} \soquete{(#4,#1)}{0}}
  \def\cm#1{{\color{unooMuted}#1}}

  % ---- componentes da plataforma
  \node[componente, text width=27mm, minimum height=44mm, anchor=north west] (hub) at (0,-4.6)
    {\estereo{componente}\\\textbf{Hub}\\{\scriptsize\itshape [Next.js]}\\[2pt]
     sessão, papel,\\\emph{embed}, assinatura\\do \emph{token}};
  \node[componente, text width=50mm, minimum height=31mm, text depth=17mm, anchor=north west] (exp) at (6.6,0)
    {\estereo{componente}\\\textbf{Experiência}\\{\scriptsize\itshape [Vite + React, \emph{bundle}]}};
  \node[componente, fill=unooSurface, text width=40mm, minimum height=10mm] (sdk) at ([yshift=7mm]exp.south)
    {\estereo{biblioteca}\\\textbf{@edu/sdk} --- ponte, HTTP, telemetria};
  \node[componente, text width=46mm, minimum height=72mm, anchor=north west] (api) at (6.6,-5.0)
    {\estereo{componente}\\\textbf{API}\\{\scriptsize\itshape [NestJS + Prisma]}\\[2pt]
     autorização por papel,\\\emph{tenancy}, regra de domínio,\\gabarito, porta da IA};

  % ---- externos
  \node[componenteext, text width=22mm, minimum height=12mm, anchor=north west] (pg) at (0,-10.3)
    {\estereo{externo}\\\textbf{PostgreSQL}\\{\scriptsize\itshape [Cloud SQL]}};
  \node[componenteext, text width=44mm, minimum height=11mm, anchor=north west] (ps) at (6.6,-14.1)
    {\estereo{externo}\\\textbf{Pub/Sub}};
  \node[componenteext, text width=22mm, minimum height=11mm, anchor=north west] (vx) at (15.3,-5.0)
    {\estereo{externo}\\\textbf{Vertex AI}};
  \node[componenteext, text width=22mm, minimum height=15mm, anchor=north west] (agp) at (15.3,-6.75)
    {\estereo{externo}\\\textbf{Agente do}\\\textbf{professor}\\{\scriptsize\itshape [ADK no}\\{\scriptsize\itshape Agent Engine]}};
  \node[componenteext, text width=22mm, minimum height=19mm, anchor=north west] (lav) at (15.3,-8.95)
    {\estereo{externo}\\\textbf{LiveAvatar}\\{\scriptsize\itshape [HeyGen]}};
  \node[componenteext, text width=22mm, minimum height=11mm, anchor=north west] (gcs) at (15.3,-11.1)
    {\estereo{externo}\\\textbf{Cloud Storage}};

  % ---- hub -> experiencia: a ponte (o hub fornece, a experiencia requer)
  \draw[associa] (1.5,0 |- hub.north) -- (1.5,-1.35) -- (4.67,-1.35);
  \draw[associa] (exp.west |- 0,-1.35) -- (5.02,-1.35);
  \bola{(4.8,-1.35)} \soquete{(4.8,-1.35)}{180}
  \node[ifacelbl, anchor=south] at (4.0,-1.15) {Ponte \texttt{postMessage} v2};
  \node[ifacelbl, anchor=north, align=left] at (4.0,-1.6)
    {\texttt{hub:init}, \texttt{hub:token},\\\texttt{hub:pause}/\texttt{resume}};

  % ---- hub <-> api (pelo vao entre os dois)
  \montD{-5.7}{3.17}{6.6}{4.9}
  \node[ifacelbl, anchor=south west, align=left] at (3.35,-5.55) {\texttt{POST /auth/login}\\\cm{\emph{Credentials}}};
  \montD{-7.2}{3.17}{6.6}{4.9}
  \node[ifacelbl, anchor=south west, align=left] at (3.35,-7.05) {REST pelo BFF\\\cm{\texttt{/api/plataforma/*}}};
  \montE{-8.7}{3.17}{6.6}{4.9}
  \node[ifacelbl, anchor=south west, align=left] at (3.35,-8.55) {JWKS \cm{\texttt{/.well-known/}}\\\cm{\texttt{jwks.json}}};

  % ---- experiencia (pela sdk) <-> api
  \montBaixo{7.7}{-3.1}{-5.0}{-4.05}
  \node[ifacelbl, anchor=east, align=right] at (7.45,-4.05) {REST com \emph{Bearer}\\\cm{rotas de domínio}};
  \montBaixo{11.1}{-3.1}{-5.0}{-4.05}
  \node[ifacelbl, anchor=east, align=right] at (10.85,-4.05) {Telemetria\\\cm{\texttt{/telemetry/batch}}\\\cm{\texttt{/telemetry/beacon}}};

  % ---- api <-> PostgreSQL e Pub/Sub
  \montE{-10.9}{2.47}{6.6}{4.9}
  \node[ifacelbl, anchor=south west, align=left] at (3.35,-10.75) {SQL + RLS\\\cm{\texttt{comTenancy}}};
  \montBaixo{7.6}{-12.2}{-14.1}{-13.15}
  \node[ifacelbl, anchor=east, align=right] at (7.35,-13.15) {Publicação\\\cm{\texttt{\{contrato\}-telemetria}}};
  \montCima{10.9}{-12.2}{-14.1}{-13.15}
  \node[ifacelbl, anchor=west, align=left] at (11.15,-13.15) {Ingestão \emph{push}\\\cm{\texttt{/telemetry/ingest}}};

  % ---- api <-> externos da direita
  \montD{-5.55}{11.67}{15.3}{13.5}
  \node[ifacelbl, anchor=south, align=left] at (13.45,-5.4) {Geração e voz ao vivo\\\cm{ADC, \texttt{:generateContent},}\\\cm{Live por WebSocket}};
  \montD{-7.5}{11.67}{15.3}{13.5}
  \node[ifacelbl, anchor=south, align=left] at (13.45,-7.35) {Turnos do agente\\\cm{\texttt{:query}, \texttt{:streamQuery}}};
  \montE{-9.45}{11.67}{15.3}{13.5}
  \node[ifacelbl, anchor=south, align=left] at (13.45,-9.3) {LLM \cm{\texttt{/ajuda/llm}}};
  \montD{-10.35}{11.67}{15.3}{13.5}
  \node[ifacelbl, anchor=north, align=left] at (13.45,-10.5) {Sessões \cm{\texttt{/v1/sessions}}};
  \montD{-11.65}{11.67}{15.3}{13.5}
  \node[ifacelbl, anchor=north, align=left] at (13.45,-11.8) {Objetos\\\cm{\emph{assets}, planos}};

\end{tikzpicture}}
\caption[Componentes e interfaces]{Diagrama de componentes: bola é interface
  fornecida, soquete é interface requerida, e o encaixe é o conector de montagem.
  A experiência requer a ponte do hub e, pela SDK, o REST e a telemetria da API;
  a API requer o JWKS do hub. O Cloud Storage também serve o \emph{bundle} ao
  navegador pela CDN, o que é artefato e não interface --- está na
  Figura~\ref{fig:uml-implantacao}. À direita, o LiveAvatar da Ajuda com avatar
  (\adrr{ADR-132}), nos dois sentidos: a API requer as sessões dele (abrir e
  fechar a sala), e ele requer o LLM da API, OpenAI-compatível. O vídeo vai do
  \emph{vendor} direto ao navegador, sem interface com a API. Também à direita,
  o agente de IA do professor no Agent Engine, que só a API chama, e o Vertex
  AI, que fornece também a voz ao vivo do mascote da Ajuda (Gemini Live,
  \adrr{ADR-158}). Fontes: \arq{hub/components/ExpFrame.tsx},
  \arq{hub/app/.well-known/jwks.json/route.ts}, \arq{api/src/auth/token.service.ts},
  \arq{api/src/telemetry/telemetry.controller.ts},
  \arq{api/src/ai/liveavatar.client.ts},
  \arq{api/src/ai/agent-engine.client.ts},
  \arq{api/src/ai/gemini-live.client.ts},
  \arq{api/src/ajuda/ajuda-llm.controller.ts} e
  \arq{packages/sdk/src/types.ts}.}
\label{fig:uml-componentes}
\end{figure}

\subsubsection{Por dentro da página da missão}

O diagrama de estrutura composta abre um classificador e mostra as
\emph{partes} que ele contém, as \emph{portas} por onde ele fala com o exterior e
os \emph{conectores} entre elas. A Figura~\ref{fig:uml-estrutura} abre o
\cd{ExpHost}, o componente cliente que a página \rota{/exp/[expId]} renderiza, e
mostra por onde cada coisa atravessa a fronteira do \emph{sandbox} (a linha
laranja). O \emph{token} entra uma vez, como propriedade vinda do servidor
(\cd{montarEmbed}), e é renovado pela porta \cd{token} do \cd{ExpFrame}, nunca
pela experiência: ela só \emph{pede}, pela ponte, e recebe \cd{hub:token} de
volta. O \emph{bundle} entra pela porta \cd{src}, direto da CDN. Entre o hub e
a experiência, é só isso que atravessa. O que a experiência fala com a API não
passa pelo hub: sai direto do \cd{iframe}, com o \emph{token} da memória
(Figura~\ref{fig:uml-componentes}).

\begin{figure}[H]
\centering
\fitwidth{%
\begin{tikzpicture}[x=1cm, y=1cm]
  \hyphenpenalty=10000 \exhyphenpenalty=10000
  \tikzset{
    rotulo/.style={font=\scriptsize, anchor=north west, inner sep=2.5pt, align=left},
    externo2/.style={rectangle, draw=unooInkSoft, line width=0.8pt, fill=unooPaper,
      font=\scriptsize, text width=30mm, inner sep=3pt, align=left},
  }

  % ---- o classificador estruturado
  \draw[draw=unooNavy, line width=1pt, fill=unooSurface] (0,0) rectangle (15.8,-9.95);
  \node[rotulo] at (0,0) {\textbf{ExpHost} \ \estereo{client component}};
  \draw[draw=unooNavy, line width=0.6pt] (0,-0.55) -- (15.8,-0.55);

  % partes
  \draw[draw=unooNavy, line width=0.85pt, fill=unooBlueWash] (0.35,-0.95) rectangle (15.45,-9.6);
  \node[rotulo] at (0.35,-0.95) {\textbf{dispositivo : DispositivoFrame [0..1]}\\
    {\color{unooMuted}só quando o manifesto declara \texttt{viewport}}};
  \draw[draw=unooNavy, line width=0.85pt, fill=unooSurface] (0.7,-2.1) rectangle (15.1,-9.25);
  \node[rotulo] at (0.7,-2.1) {\textbf{frame : ExpFrame [1]}};

  % o iframe: a fronteira do sandbox
  \draw[iframe] (1.1,-2.95) rectangle (8.9,-8.85);
  \node[iframelbl] at (1.2,-2.8) {\texttt{iframe [1]}};
  \node[font=\scriptsize, anchor=north west, text width=72mm, align=left] at (1.25,-3.3)
    {{\color{unooMuted}um \texttt{HTMLIFrameElement} com}\\
     \texttt{sandbox = "allow-scripts allow-same-origin}\\
     \texttt{\ \ allow-modals"} {\color{unooMuted}+ \texttt{pointer-lock}, \texttt{forms},}\\
     {\color{unooMuted}\ \ \texttt{downloads} se o manifesto pedir}\\
     \texttt{allow = "autoplay; fullscreen; clipboard-write"}};
  \node[parte, text width=40mm, anchor=north west, fill=unooGreenWash, draw=unooGreen!60!black]
      (bexp) at (1.5,-5.85)
    {\textbf{exp : \estereo{Experiencia}}\\{\color{unooMuted}o \emph{bundle}, servido pela CDN}};
  \node[parte, text width=24mm, anchor=north west] (bsdk) at (1.5,-7.45)
    {\textbf{sdk : @edu/sdk}};
  \node[parte, text width=22mm, anchor=north west] (btok) at (4.95,-7.45)
    {\textbf{token : string}\\{\color{unooMuted}\{só em memória\}}};

  % partes do frame fora do sandbox
  \node[parte, text width=42mm, anchor=north west, fill=unooOrangeWash, draw=unooOrange] (aviso) at (10.0,-2.95)
    {\textbf{aviso : AlertBox [0..1]}\\{\color{unooMuted}\enquote{Sua sessão expirou};}\\{\color{unooMuted}substitui o \texttt{iframe}}};
  \node[parte, text width=42mm, anchor=north west] (proto) at (10.0,-5.85)
    {\textbf{ponte : protocolo v2}\\{\color{unooMuted}\texttt{postMessage}; \texttt{ev.origin} + Zod}};

  % ---- portas
  \node[porta] (pEmbed) at (1.8,0) {};
  \node[porta] (pProg) at (14.0,0) {};
  \node[porta] (pTok) at (10.6,-9.25) {};
  \node[porta] (pLog) at (14.1,-9.25) {};
  \node[portaborda] (pPonte) at (8.9,0 |- proto.west) {};
  \node[portaborda] (pSrc) at (8.0,-2.95) {};
  \node[font=\scriptsize\bfseries\color{unooOrange}, anchor=south east, align=right] at (pPonte.north west) {ponte};
  \node[font=\scriptsize\bfseries\color{unooOrange}, anchor=north west] at (pSrc.south east) {src};

  % ---- externos (em cima: quem entrega; embaixo: as rotas do hub)
  \node[externo2, text width=34mm, anchor=south west] (page) at (0,0.85)
    {\estereo{server component}\\\textbf{app/exp/[expId]/page.tsx}\\\texttt{montarEmbed()}: \emph{token},\\\texttt{src}, \texttt{expOrigin},\\\texttt{permissoes}, \texttt{viewport}};
  \node[externo2, text width=37mm, anchor=south] (cdn) at (8.0,0.85)
    {\estereo{CDN}\\\texttt{\{cdn\}/\{expId\}/\{versao\}/}\\\texttt{index.html?hubOrigin=\ldots}};
  \node[externo2, text width=31mm, anchor=south] (rprog) at (14.0,0.85)
    {\estereo{BFF}\\\textbf{POST /api/plataforma/}\\\textbf{progresso/:expId/}\\\textbf{concluir}};
  \node[externo2, text width=31mm, anchor=north] (rtok) at (10.6,-10.85)
    {\estereo{route handler}\\\textbf{GET /api/exp-token}\\401/403 \seta{} \texttt{hub:pause}\\e o aviso};
  \node[externo2, text width=31mm, anchor=north] (rlog) at (14.1,-10.85)
    {\estereo{route handler}\\\textbf{POST /api/ponte}\\uma linha de log; 204};

  % ---- conectores
  \draw[conector] (page.south -| pEmbed) -- (pEmbed.north) node[cxlabel, midway, right=2pt]{embed};
  \draw[conector] (rprog.south) -- (pProg.north) node[cxlabel, midway, right=2pt]{progresso};
  \draw[conector] (pTok.south) -- (rtok.north) node[cxlabel, midway, right=2pt]{token};
  \draw[conector] (pLog.south) -- (rlog.north) node[cxlabel, midway, right=2pt]{sinal};
  \draw[conector, draw=unooOrange, line width=1pt] (pPonte.east) -- (proto.west);
  \draw[conector] (pTok.north) -- (pTok.north |- proto.south);
  \draw[conector] (pLog.north) -- (pLog.north |- proto.south);
  \draw[conector, draw=unooOrange, line width=1pt] (cdn.south) -- (pSrc.north);
  \draw[conector] (pPonte.west) -- (bexp.east |- pPonte);
  \draw[conector] (bexp.south -| bsdk.north) -- (bsdk.north);
  \draw[conector] (bsdk.east) -- (btok.west);

  \node[actnote, text width=37mm, anchor=north west] at (0,-10.4)
    {\textbf{\texttt{targetOrigin} é sempre \texttt{expOrigin}}, nunca \texttt{*}: o
     \texttt{frame} só envia para a origem que \texttt{montarEmbed} calculou, e só
     aceita mensagem cujo \texttt{ev.origin} seja ela.};
  \node[actnote, text width=40mm, anchor=north west] at (4.2,-10.4)
    {\textbf{O \texttt{ExpHost} tem uma segunda sobreposição}, fora do
     desenho: \enquote{A missão travou}, com o \emph{hash} do erro, quando a
     experiência manda \texttt{exp:erro}. É dele, não do \texttt{frame}: sobrevive
     à remontagem do \texttt{frame} ao recarregar.};

\end{tikzpicture}}
\caption[Estrutura composta: a página da missão]{Diagrama de estrutura composta do
  \cd{ExpHost}. Partes são instâncias dentro do classificador (com multiplicidade);
  quadrados na borda são portas; a linha laranja tracejada é a fronteira do
  \emph{sandbox}, e as portas laranja são as duas únicas ligações entre o hub e
  o que roda dentro dela; as chamadas da experiência à API saem direto do
  \texttt{iframe}. Fontes: \arq{hub/components/ExpHost.tsx},
  \arq{hub/components/ExpFrame.tsx}, \arq{hub/components/DispositivoFrame.tsx},
  \arq{hub/lib/embed-experiencia.ts} e \arq{hub/app/api/exp-token/route.ts}.}
\label{fig:uml-estrutura}
\end{figure}

% ---------------------------------------------------------------------------
\subsection{Visão de processos: atividades, interações e tempo}
\label{sec:uml-processos}

\subsubsection{O algoritmo da disponibilidade}

O diagrama de atividades da Figura~\ref{fig:uml-atividade-disp} é o
\cd{exigirRecursoDisponivel}, a primeira das duas travas de toda entrada em
experiência. As duas \emph{partições} são os dois arquivos, e é a divisão que
importa: a da esquerda lê o banco dentro da transação de quem chama; a da direita
é função pura, testada em milissegundos sem PostgreSQL. As regras da
Tabela~\ref{tab:uml-objetos} moram inteiras na partição da direita.

\begin{figure}[H]
\centering
\fitpage{%
\begin{tikzpicture}[x=1cm, y=1cm]
  \hyphenpenalty=10000 \exhyphenpenalty=10000
  \tikzset{
    ac/.style={action, text width=40mm},
    dc/.style={decision, text width=22mm, aspect=2.2},
    flowfinal/.style={circle, draw=unooInk, line width=0.9pt, fill=unooSurface,
      minimum size=4.6mm, inner sep=0pt,
      path picture={\draw[unooInk, line width=0.8pt]
        (path picture bounding box.south west) -- (path picture bounding box.north east)
        (path picture bounding box.north west) -- (path picture bounding box.south east);}},
  }

  % ---- particoes
  \draw[lane] (0,0.2) rectangle (7.4,-21.2);
  \draw[lane] (7.4,0.2) rectangle (18.0,-21.2);
  \node[lanehdr, minimum width=74mm, minimum height=7mm, anchor=north west] at (0,0.2)
    {\texttt{comum/disponibilidade.ts} --- com banco};
  \node[lanehdr, minimum width=106mm, minimum height=7mm, anchor=north west] at (7.4,0.2)
    {\texttt{recursos/regras.ts} --- função pura};

  % ---- particao da esquerda: ler
  \node[actstart] (ini) at (3.7,-1.2) {};
  \node[ac] (a1) at (3.7,-2.3) {buscar \texttt{experiencias[expId]}\\{\color{unooMuted}\texttt{produto}, \texttt{publico}; pode não existir}};
  \node[ac] (a2) at (3.7,-3.8) {\texttt{cadeiaDaEscola}\\{\color{unooMuted}\texttt{[escola]} ou \texttt{[escola, regional]}}};
  \node[ac] (a3) at (3.7,-5.3) {\texttt{linhasDoEscopo}\\{\color{unooMuted}escola nula \textbf{ou} na cadeia; só \emph{ids} legíveis}};
  \node[ac] (a4) at (3.7,-6.8) {\texttt{produtoConhecido}\\{\color{unooMuted}fora do vocabulário \seta{} \texttt{null}}};

  % ---- particao da direita: decidir
  \node[dc] (d1) at (12.4,-7.8) {produto?};
  \node[ac] (a5) at (12.4,-9.4) {\texttt{maisEspecifica(}\\\texttt{"produto:<slug>")}\\{\color{unooMuted}sem linha: \texttt{ligadoPorPadrao}}};
  \node[dc] (d2) at (12.4,-11.0) {produto ligado?};
  \node[actionerr, text width=25mm] (r1) at (16.45,-11.0) {desligado,\\\texttt{tetoDoProduto}};
  \node[escolha] (m1) at (12.4,-12.5) {};
  \node[ac] (a6) at (12.4,-13.6) {\texttt{maisEspecifica(}\\\texttt{"experiencia:<expId>")}};
  \node[dc] (d3) at (12.4,-15.2) {linha da\\experiência?};
  \node[actionok, text width=26mm] (r2) at (9.4,-16.8) {ligado,\\\texttt{versaoFixa = null}};
  \node[action, text width=32mm] (r3) at (13.4,-16.8) {\texttt{ligado} e \texttt{versaoFixa}\\da linha};
  \node[escolha] (m2) at (12.4,-18.4) {};

  % ---- de volta a esquerda: agir
  \node[dc] (d4) at (3.7,-19.1) {\texttt{efetivo.ligado}?};
  \node[actend] (fim) at (0.9,-19.1) {};
  \node[actionerr, text width=40mm] (a7) at (3.7,-20.5) {lança 403 \texttt{recurso-indisponivel}};
  \node[flowfinal] (ff) at (6.6,-20.5) {};

  % ---- fluxos
  \draw[flow] (ini) -- (a1);
  \draw[flow] (a1) -- (a2);
  \draw[flow] (a2) -- (a3);
  \draw[flow] (a3) -- (a4);
  \draw[flow] (a4.east) -- (12.4,-6.8) -- (d1.north);
  \node[flowlbl, anchor=south] at (9.3,-6.8) {\texttt{resolverExperiencia}};
  \draw[flow] (d1) -- (a5) node[flowlbl, midway, right=2pt]{[sim]};
  \draw[flow] (a5) -- (d2);
  \draw[flow] (d2) -- (r1) node[flowlbl, midway, above]{[não]};
  \draw[flow] (d2) -- (m1) node[flowlbl, midway, right=2pt]{[sim]};
  \draw[flow] (d1.west) -- node[flowlbl, above]{[não]} (9.2,-7.8) -- (9.2,-12.5) -- (m1.west);
  \draw[flow] (m1) -- (a6);
  \draw[flow] (a6) -- (d3);
  \draw[flow] (d3.west) -| (r2.north) node[flowlbl, pos=0.3, above]{[nenhuma]};
  \draw[flow] (d3.south) -- (d3.south |- r3.north) node[flowlbl, midway, right=2pt]{[uma linha]};
  \draw[flow] (r2.south) |- (m2.west);
  \draw[flow] (r3.south -| 12.4,0) -- (m2.north);
  \draw[flow] (r1.south) -- (16.45,-18.4) -- (m2.east);
  \draw[flow] (m2.south) |- (d4.east);
  \draw[flow] (d4) -- (fim) node[flowlbl, midway, above]{[sim]};
  \draw[flow] (d4) -- (a7) node[flowlbl, midway, left=2pt]{[não]};
  \draw[flow] (a7) -- (ff);

  % ---- notas
  \node[actnote, text width=46mm, anchor=north west] (n1) at (12.9,-0.9)
    {\textbf{\texttt{maisEspecifica}} percorre a cadeia na ordem --- escola,
     depois regional --- e devolve a primeira linha do recurso; sem nenhuma,
     a de escola nula (o contrato); sem ela, \texttt{null}. Sobe só um nível:
     a rede não tem árvore mais funda.};
  \node[actnote, text width=40mm, anchor=north west] (n2) at (7.8,-0.9)
    {\textbf{A linha inteira é a decisão:} \texttt{ligado} e \texttt{versaoFixa}
     não se resolvem em cadeias separadas. A tela mostra exatamente a linha que
     vale.};
  \node[actnote, text width=30mm, anchor=north west] at (0.35,-8.2)
    {\textbf{Sem linha em \texttt{experiencias}} (\emph{stack} nova, antes do
     primeiro \emph{publish}) \textbf{não tranca}: sem produto, vale só o
     interruptor da própria experiência.};
  \node[actnote, text width=30mm, anchor=north west] at (3.9,-8.2)
    {\textbf{A mensagem depende do público:} \enquote{Este recurso não está
     disponível para a sua rede.} para professor; \enquote{Esta missão não está
     aberta agora.} para estudante.};

\end{tikzpicture}}
\caption[Atividades: o resolvedor de disponibilidade]{Diagrama de atividades do
  \cd{exigirRecursoDisponivel}, com uma partição por arquivo. A da direita é
  \cd{resolverExperiencia}, pura; os losangos vazios são nós de junção. A
  terminação em \emph{X} é o fim de fluxo pela exceção: a transação de quem
  chamou desfaz o que tiver feito. O produto sem linha vale o padrão dele
  (\cd{ligadoPorPadrao}): ligado, salvo os de
  \cd{PRODUTOS\_\B DESLIGADOS\_\B POR\_\B PADRAO} (hoje só a \cd{ajuda},
  \adrr{ADR-132}); a experiência sem linha continua ligada. Fontes:
  \arq{comum/disponibilidade.ts}, \arq{recursos/regras.ts} e
  \arq{recursos/catalogo.ts}.}
\label{fig:uml-atividade-disp}
\end{figure}

\begin{nota}
\textbf{O produto sem experiência entra por um atalho do mesmo algoritmo.} O
agente de IA do professor e a Ajuda com avatar não têm linha em
\cd{experiencias}, então a classe declara \cd{@Produto("agente-professor")} ou
\cd{@Produto("ajuda")} em vez de \cd{@Recurso(expId)}, e o
\cd{DisponibilidadeGuard} chama \cd{exigirProdutoDisponivel}: só o ramo do
TETO do diagrama acima (\cd{disponibilidadeDoProduto} \seta{} \cd{resolverProduto},
escola \seta{} regional \seta{} contrato \seta{} padrão do produto), sem a etapa
da experiência. Para a Ajuda o padrão é \emph{desligado} (\adrr{ADR-132}), e o
\emph{endpoint} que o \emph{vendor} chama refaz a mesma pergunta a cada turno, pela
escola do \emph{token} da sessão: desligar no meio da conversa vira uma fala do
atendente, não um erro. É o MESMO resolvedor de propósito: uma
segunda leitura da tabela com regra própria seria o painel dizendo
\enquote{desligado} e a porta abrindo. Declarar os dois decoradores na mesma
classe lança, como declarar nenhum.
\end{nota}

\subsubsection{A entrada na missão: onde cada trava dispara}

A regra 7 do \arq{AGENTS.md} diz que filtrar o \emph{card} no hub é UX e que a
trava é da API. O diagrama de sequência da Figura~\ref{fig:uml-seq-entrada} mostra
por quê, na ordem em que as coisas acontecem. O hub consulta a grade e devolve 404
para missão que não está liberada --- mas é a mesma grade que o estudante vê, e o
\emph{bundle} é público: quem digita \rota{/exp/:expId} de uma missão fechada, ou
chama a API direto, só encontra a trava de verdade na primeira rota de estudante,
no funil \cd{garantirSessao}/\cd{estudanteDaEscola}. Ali as duas travas disparam
em série, disponibilidade antes de liberação.

\begin{figure}[H]
\centering
\fitpage{%
\begin{tikzpicture}[x=1cm, y=1cm]

  \tikzset{guarda/.append style={align=left}}

  \raia{1.3}{-22.9}{Navegador}{hub no cliente}
  \raia{5.0}{-22.9}{Hub}{servidor}
  \raia{12.2}{-22.9}{API}{NestJS}
  \raiaext{16.4}{-22.9}{PostgreSQL}{Cloud SQL}

  % ---- a grade
  \msgs{1.3}{5.0}{-1.1}{\texttt{GET /exp/:expId} (\emph{cookie})}
  \msgs{5.0}{12.2}{-1.9}{\texttt{GET /salas/minhas}}
  \msgs{12.2}{16.4}{-2.9}{\texttt{turmasDoEstudante}\\+ \texttt{liberacoes}}
  \msgr{16.4}{12.2}{-3.6}{matrículas e linhas}
  \msgr{12.2}{5.0}{-4.4}{grade: estado por missão, sem as indisponíveis}

  % ---- alt: o filtro do hub
  \draw[seqframe] (0.4,-4.85) rectangle (13.1,-10.2);
  \node[frameop] at (0.4,-4.85) {alt};
  \node[guarda] at (1.45,-5.25) {[missão não está\\\texttt{liberada} na grade]};
  \msgr{5.0}{1.3}{-6.2}{404 (\texttt{notFound})}
  \draw[seqframe, dashed] (0.4,-6.5) -- (13.1,-6.5);
  \node[guarda] at (1.45,-6.77) {[\texttt{liberada}]};
  \msgself{5.0}{-7.15}{\texttt{montarEmbed}: assina o \emph{token} (RS256, 15 min)}
  \msgs{5.0}{12.2}{-8.25}{\texttt{GET /experiencias/:expId} e \texttt{/recursos/efetivos}}
  \msgr{12.2}{5.0}{-9.05}{manifesto; versão fixada da escola}
  \msgr{5.0}{1.3}{-9.85}{HTML com \texttt{ExpFrame}}

  % ---- a experiencia nasce
  \node[seqheadact] (hexp) at (8.5,-11.1) {\textbf{Experiência}\\[0.5pt]{\scriptsize\itshape \texttt{iframe}}};
  \draw[seqline] (8.5,-11.62) -- (8.5,-22.9);
  \draw[seqret] (1.3,-11.1) -- (hexp.west)
    node[seqlbl, pos=0.36, above=0.4pt]{\estereo{create} \emph{bundle} da CDN};

  % ---- ref: o handshake
  \draw[seqframe, fill=unooSurface] (0.4,-12.0) rectangle (9.4,-12.9);
  \node[frameop] at (0.4,-12.0) {ref};
  \node[font=\scriptsize] at (5.2,-12.47)
    {\emph{handshake} \texttt{exp:ready} / \texttt{hub:init} --- Figura~\ref{fig:seq-handshake}};

  % ---- a primeira rota de estudante: as duas travas
  \msgs{8.5}{12.2}{-13.7}{1ª rota de estudante (\emph{Bearer})}
  \msgself{12.2}{-14.3}{funil: \texttt{garantirSessao}\\ou \texttt{estudanteDaEscola}}
  \msgs{12.2}{16.4}{-16.05}{\texttt{exigirRecursoDisponivel}\\(Figura~\ref{fig:uml-atividade-disp})}
  \msgs{12.2}{16.4}{-16.85}{\texttt{missoesLiberadasPara}}
  \msgr{16.4}{12.2}{-17.85}{linhas de disponibilidade\\e de liberação}

  \draw[seqframe] (7.6,-18.25) rectangle (17.4,-22.6);
  \node[frameop] at (7.6,-18.25) {alt};
  \node[guarda] at (8.65,-18.65) {[recurso desligado no\\escopo que vale]};
  \msgr{12.2}{8.5}{-19.6}{403 \texttt{recurso-indisponivel}}
  \draw[seqframe, dashed] (7.6,-19.9) -- (17.4,-19.9);
  \node[guarda] at (8.65,-20.28) {[sem liberação \texttt{liberada}\\nas turmas do estudante]};
  \msgr{12.2}{8.5}{-21.15}{403 \texttt{missao-nao-liberada}}
  \draw[seqframe, dashed] (7.6,-21.45) -- (17.4,-21.45);
  \node[guarda] at (8.65,-21.72) {[as duas passam]};
  \msgr{12.2}{8.5}{-22.3}{200 --- a rota segue}

\end{tikzpicture}}
\caption[Sequência: a entrada do estudante na missão]{Sequência da entrada do
  estudante na missão. O filtro do hub (o primeiro \texttt{alt}) é UX; as travas
  estão no segundo, na API, e disparam em série. A linha de vida da experiência
  começa na mensagem \estereo{create}: antes dela, o \texttt{iframe} não existe.
  Fontes: \arq{hub/app/exp/[expId]/page.tsx}, \arq{hub/lib/embed-experiencia.ts},
  \arq{api/src/comum/liberacao.ts}, \arq{api/src/comum/disponibilidade.ts} e
  \arq{api/src/salas/salas.service.ts}.}
\label{fig:uml-seq-entrada}
\end{figure}

\begin{nota}
\textbf{Os dois 403 chegam iguais à experiência.} A SDK não repassa o corpo do
erro: \cd{call} lança \cd{Error("API GET \ldots{} \seta{} 403")}, e cada experiência tira
o \emph{status} da mensagem. A exp05ef1b e a exp04ef1b mostram o aviso dentro do
disfarce (\enquote{Este grupo ainda não abriu para a sua turma}); o hub não mostra
nada, porque não existe mensagem de ponte para isso. Se um dia a diferença entre
\enquote{não liberada} e \enquote{indisponível na rede} precisar chegar ao
estudante, o caminho é a SDK expor o \cd{codigo}, não a experiência adivinhar.
\end{nota}

\subsubsection{Uma chamada de IA: o que roda antes e depois do \emph{provider}}

A Figura~\ref{fig:uml-seq-ia} é a única porta de IA aberta ao cliente,
\rota{POST /ai/gerar}, na ordem do código. Os dois fragmentos \texttt{break} são
as saídas antecipadas: quando a guarda do \texttt{break} vale, o operando roda e o
resto da interação \emph{não} acontece. O que o diagrama deixa ver e o texto
esconde: a leitura do teto (\cd{ia\_quotas}) roda \emph{fora} do
\cd{comTenancy}, porque a linha é do contrato e não tem escola; a do consumo
(\cd{ia\_uso\_diario}) roda dentro, porque tem. E o prompt de sistema é escolhido
e montado dentro do \emph{controller}, com a série resolvida da turma do
\emph{token} --- nunca vinda do corpo.

\begin{figure}[H]
\centering
\fitpage{%
\begin{tikzpicture}[x=1cm, y=1cm]

  \tikzset{seqhead/.append style={text width=22mm},
           guarda/.append style={align=left}}
  % auto-mensagem estreita: o rotulo cabe entre esta linha de vida e a vizinha
  \def\selfc#1#2#3{%
    \draw[seqself] (#1,#2) -- ++(0.45,0) -- ++(0,-0.4) -- ++(-0.45,0);
    \node[seqlbl, anchor=west, align=left] at ({#1+0.52},{#2-0.2}) {#3};}
  \def\fundo{-24.15}

  \raia{1.25}{\fundo}{Experiência}{pela SDK}
  \raia{3.75}{\fundo}{\texttt{OidcGuard}}{\texttt{auth/}}
  \raia{6.7}{\fundo}{\texttt{AiController}}{\texttt{ai/}}
  \raia{9.35}{\fundo}{\texttt{QuotaService}}{\texttt{ai/}}
  \raiaext{12.45}{\fundo}{PostgreSQL}{Cloud SQL}
  \node[seqhead, text width=26mm] at (15.15,0) {\textbf{\texttt{provider}}\\[0.5pt]{\scriptsize\itshape \texttt{MedidorDeProvider}}};
  \draw[seqline] (15.15,-0.52) -- (15.15,\fundo);
  \raiaext{18.1}{\fundo}{Vertex AI}{Gemini}

  \msgs{1.25}{3.75}{-1.3}{\texttt{POST /ai/gerar}\\(\emph{Bearer})}
  \selfc{3.75}{-1.75}{\texttt{verificar}\\(JWKS do hub)}

  \draw[seqframe] (0.4,-2.65) rectangle (7.6,-4.15);
  \node[frameop] at (0.4,-2.65) {break};
  \node[guarda] at (1.55,-3.03) {[\emph{token} ausente\\ou inválido]};
  \msgr{3.75}{1.25}{-3.9}{401}

  \msgs{3.75}{6.7}{-5.35}{\texttt{req.auth}:\\usuário, escola\\e perfil}
  \selfc{6.7}{-5.75}{\texttt{GerarSchema}\\\texttt{.parse} (Zod;\\400 se inválido)}
  \msgs{6.7}{9.35}{-7.55}{\texttt{temSaldo}(escola,\\usuário, perfil)}
  \msgs{9.35}{12.45}{-8.45}{\texttt{limite}: \texttt{ia\_quotas}\\(fora do \texttt{comTenancy})}
  \msgr{12.45}{9.35}{-9.1}{\texttt{tokens\_por\_dia}}
  \msgs{9.35}{12.45}{-10.5}{\texttt{usadoHoje}:\\\texttt{ia\_uso\_diario}\\(em \texttt{comTenancy})}
  \msgr{12.45}{9.35}{-11.15}{\texttt{tokens\_usados}}
  \msgr{9.35}{6.7}{-11.8}{\texttt{usado < limite}}

  \draw[seqframe] (0.4,-12.15) rectangle (10.3,-13.55);
  \node[frameop] at (0.4,-12.15) {break};
  \node[guarda] at (1.55,-12.42) {[sem saldo]};
  \draw[seqret] (6.7,-13.3) -- (1.25,-13.3)
    node[seqlbl, pos=0.24, above=0.4pt]{429 “Quota diária\\de IA atingida”};

  \selfc{6.7}{-13.9}{\texttt{promptDoTutor}\\\texttt{(expId)}}
  \draw[seqmsg] (6.7,-15.5) -- (12.45,-15.5)
    node[seqlbl, pos=0.23, above=0.4pt]{\texttt{serieDoPedido}:\\\texttt{turma.findFirst}};
  \draw[seqret] (12.45,-16.15) -- (6.7,-16.15)
    node[seqlbl, pos=0.77, above=0.4pt]{\texttt{serie} \seta{} \enquote{6º ano}};
  \draw[seqmsg] (6.7,-17.15) -- (15.15,-17.15)
    node[seqlbl, pos=0.157, above=0.4pt]{\texttt{generate}\\(\texttt{system}, \texttt{user})};
  \selfc{15.15}{-17.5}{\texttt{medir}:\\\texttt{VertexProvider}}
  \msgs{15.15}{18.1}{-19.2}{\texttt{:generateContent}\\(ADC)}
  \msgr{18.1}{15.15}{-19.85}{texto e uso}
  \draw[seqret] (15.15,-20.8) -- (6.7,-20.8)
    node[seqlbl, pos=0.5, above=0.4pt]{\texttt{AiResult}:\\texto e \texttt{tokensUsados}};
  \msgs{6.7}{9.35}{-21.75}{\texttt{consumir}(\ldots,\\\texttt{tokensUsados})}
  \msgs{9.35}{12.45}{-22.45}{\emph{upsert} em \texttt{comTenancy}}
  \msgr{12.45}{9.35}{-23.1}{\texttt{tokens\_usados}}
  \draw[seqret] (6.7,-23.8) -- (1.25,-23.8)
    node[seqlbl, pos=0.24, above=0.4pt]{200 \texttt{\{texto\}}};

\end{tikzpicture}}
\caption[Sequência: uma chamada de IA com quota]{Sequência de \rota{POST /ai/gerar}.
  \texttt{break} é a saída antecipada: 401 sem \emph{token} válido, 429 sem saldo.
  O consumo é gravado \emph{depois} da resposta do modelo, com os \emph{tokens} que
  ele declarou; entre o \texttt{temSaldo} e o \texttt{consumir} não há trava. A
  Ajuda com avatar é a única IA que não passa por este desenho: é limitada por
  minutos, não por \emph{tokens} (\adrr{ADR-132},
  Figura~\ref{fig:uml-seq-ajuda}). Fontes:
  \arq{ai/ai.controller.ts}, \arq{ai/prompts.ts}, \arq{ai/quota.service.ts},
  \arq{ai/medidor.provider.ts}, \arq{auth/oidc.guard.ts}.}
\label{fig:uml-seq-ia}
\end{figure}

\begin{limitacao}
\textbf{\rota{/ai/gerar} não tem \emph{rate limit}.} A quota limita o custo do dia;
a rajada dentro do dia só é contida nas rotas que chamam \cd{rateLimit} ---
\arq{planos}, \arq{bemestar}, \arq{manifesto-ia}, \arq{vila} e as reações do
\arq{motor}. O \arq{unoogpt} chama \cd{rateLimit} e ignora o retorno, o que faz o
limite de 12 por minuto que o contrato da rota promete não valer. E o
\cd{Map} do \cd{rateLimit} é por instância: com \cd{min-instances} acima de um, o
limite efetivo é multiplicado. Os três pontos estão fora desta frente; ficam
registrados aqui para quem desenhar a próxima rota de IA.
\end{limitacao}

\subsubsection{Um turno do agente de IA do professor: a quota antes e depois do \emph{stream}}

A Figura~\ref{fig:uml-seq-agente} é \rota{POST /agente-professor/conversas/:id/mensagens},
a única rota da plataforma que responde \cd{text/event-stream} (\adrr{ADR-128}),
na ordem do código. Ela é a Figura~\ref{fig:uml-seq-ia} com três diferenças que só
o diagrama deixa ver juntas. \textbf{Toda recusa sai antes do \emph{stream}}, em
JSON (401 sem \emph{token}; 403 sem o papel \cd{professor} ou com o produto
desligado; 400 corpo inválido; 404 conversa de outro professor ou de outra
escola; 429 acima de 10 mensagens por minuto, \cd{RATE\_LIMIT}, ou sem saldo,
\cd{QUOTA}): depois do 200 só existe o evento \cd{erro}. \enquote{Parar} (ou
fechar a aba) aborta cada salto até o Vertex; fechar o \emph{painel} não aborta
--- o turno segue e acende o selo do botão. \textbf{O consumo da quota mora
num \cd{finally}}, e soma os \emph{tokens} dos eventos NÃO parciais --- o professor
pode apertar \enquote{Parar} ou fechar a aba no meio, e o que o modelo gerou foi
pago. E \textbf{a posse é
conferida no Postgres}, porque o Agent Engine não confere o usuário numa sessão
pedida por nome (\adrr{ADR-126}). O laço de ferramentas dentro do agente aparece
com uma ferramenta só; num turno real o modelo pode chamar \cd{buscar\_habilidades}
e \cd{buscar\_acervo} mais de uma vez.

\begin{figure}[H]
\centering
\fitpage{%
\begin{tikzpicture}[x=1cm, y=1cm]

  \tikzset{seqhead/.append style={text width=22mm},
           guarda/.append style={align=left}}
  \def\fundo{-24.55}
  % mensagem com o rotulo numa posicao escolhida (fora das linhas de vida que ela cruza)
  \def\msgp#1#2#3#4#5{\draw[seqmsg] (#1,#3) -- (#2,#3) node[seqlbl, pos=#4, above=0.4pt]{#5};}
  \def\retp#1#2#3#4#5{\draw[seqret] (#1,#3) -- (#2,#3) node[seqlbl, pos=#4, above=0.4pt]{#5};}

  \raia{1.25}{\fundo}{Navegador}{painel do chat}
  \raia{3.75}{\fundo}{Hub}{rota de \emph{stream}}
  \node[seqhead, text width=24mm] at (6.6,0) {\textbf{\texttt{AgenteProfessor}}\\[0.5pt]{\scriptsize\itshape \emph{controller} + \emph{service}}};
  \draw[seqline] (6.6,-0.52) -- (6.6,\fundo);
  \raiaext{9.4}{\fundo}{PostgreSQL}{Cloud SQL}
  \node[seqhead, text width=27.5mm] at (12.7,0) {\textbf{Agente}\\[0.5pt]{\scriptsize\itshape ADK no Agent Engine}};
  \draw[seqline] (12.7,-0.52) -- (12.7,\fundo);
  \raiaext{15.7}{\fundo}{Gemini}{\texttt{global} \seta{} \texttt{us}}
  \raiaext{18.2}{\fundo}{\emph{Data stores}}{Vertex AI Search}

  % ---- tudo o que pode recusar, antes do stream
  \msgs{1.25}{3.75}{-1.1}{\texttt{POST} \{\texttt{texto}\}}
  \msgs{3.75}{6.6}{-2.05}{\texttt{POST} \ldots\texttt{/mensagens}\\(\emph{token} curto)}
  \msgs{6.6}{9.4}{-3}{\texttt{daPosse}\\(dono e escola)}
  \msgr{9.4}{6.6}{-3.7}{\texttt{sessao\_id}}
  \msgs{6.6}{9.4}{-4.65}{\texttt{temSaldo}\\(perfil \texttt{professor})}
  \msgr{9.4}{6.6}{-5.35}{\texttt{usado < limite}}

  \draw[seqframe] (0.4,-5.75) rectangle (7.5,-7.15);
  \node[frameop] at (0.4,-5.75) {break};
  \node[guarda] at (1.55,-6.02) {[alguma recusa]};
  \retp{6.6}{1.25}{-6.9}{0.24}{4xx em JSON,\\antes do \emph{stream}}

  % ---- o stream
  \retp{6.6}{1.25}{-8.2}{0.265}{200 \texttt{event-stream}\\(repassado pelo hub)}
  \msgp{6.6}{12.7}{-9.05}{0.73}{\texttt{:streamQuery?alt=sse}\\\texttt{user\_id}, \texttt{session\_id}}
  \msgself{12.7}{-9.55}{guarda de escopo\\(modelo leve)}
  \msgs{12.7}{15.7}{-10.8}{gerar, em \emph{stream}}
  \msgr{15.7}{12.7}{-11.5}{\texttt{function\_call}}
  \msgp{12.7}{18.2}{-12.55}{0.27}{\texttt{buscar\_habilidades}\\ou \texttt{buscar\_acervo}}
  \retp{18.2}{12.7}{-13.45}{0.73}{habilidades\\ou trechos}
  \msgs{12.7}{15.7}{-14.15}{com o retorno}
  \msgr{15.7}{12.7}{-14.85}{texto em pedaços}

  \draw[seqframe] (0.4,-15.25) rectangle (13.6,-17.55);
  \node[frameop] at (0.4,-15.25) {loop};
  \node[guarda] at (1.55,-15.65) {[cada evento\\do agente]};
  \retp{12.7}{6.6}{-16.4}{0.27}{evento: parcial, final\\ou \texttt{state\_delta}}
  \retp{6.6}{1.25}{-17.25}{0.265}{\texttt{event: delta}\\(sem repetir o final)}

  \draw[seqframe] (0.4,-17.95) rectangle (10.3,-22.1);
  \node[frameop] at (0.4,-17.95) {alt};
  \node[guarda] at (1.45,-18.35) {[o agente\\terminou]};
  \retp{6.6}{1.25}{-19}{0.24}{\texttt{event: fontes}}
  \msgs{6.6}{9.4}{-19.9}{\texttt{atualizada\_em},\\\texttt{titulo}}
  \retp{6.6}{1.25}{-20.5}{0.24}{\texttt{titulo}, \texttt{aviso}, \texttt{fim}}
  \draw[seqframe, dashed] (0.4,-20.8) -- (10.3,-20.8);
  \node[guarda] at (1.45,-21.2) {[falha do agente\\ou 240\,s]};
  \retp{6.6}{1.25}{-21.8}{0.24}{\texttt{event: erro}}

  \msgs{6.6}{9.4}{-23.35}{\texttt{finally}: \texttt{consumir}\\(soma dos\\não parciais)}
  \msgp{6.6}{12.7}{-24.25}{0.73}{\texttt{PATCH expireTime}\\(só depois de \texttt{fim})}

\end{tikzpicture}}
\caption[Sequência: um turno do agente de IA do professor]{Sequência de
  \rota{POST /agente-professor/conversas/:id/mensagens}. \texttt{break} é a saída
  antecipada, sempre em JSON e antes do \emph{stream}; depois do 200 só há
  eventos, e a falha vira \texttt{event: erro} (\texttt{AGENTE\_INDISPONIVEL} ou,
  no teto de 240\,s, \texttt{TEMPO\_ESGOTADO}). No \texttt{loop} a api escreve
  \texttt{: ping} a cada 15\,s. O
  consumo roda no \texttt{finally} também no abandono e no erro; a renovação só
  depois de \texttt{fim}, e a falha dela é só log. Nenhuma transação
  \texttt{comTenancy} fica aberta durante o \emph{stream}. Fontes:
  \arq{agente-professor/agente-professor.controller.ts},
  \arq{agente-professor/agente-professor.service.ts},
  \arq{agente-professor/regras.ts}, \arq{ai/agent-engine.client.ts},
  \arq{apps/agente-professor/agente\_professor/plugins.py}.}
\label{fig:uml-seq-agente}
\end{figure}

\subsubsection{Uma sessão da Ajuda com avatar: minutos em vez de \emph{tokens}, e um LLM que o \emph{vendor} chama}

A Figura~\ref{fig:uml-seq-ajuda} é a Ajuda com avatar (\adrr{ADR-132}) do clique
em \enquote{Ajuda} ao fechar do modal, na ordem do código. Ela inverte duas coisas
das duas figuras anteriores. \textbf{Quem chama o modelo não é a pessoa, é o
\emph{vendor}}: o LiveAvatar faz a escuta, a voz e o vídeo (FULL Mode), e a cada fala
chama de volta \rota{POST /ajuda/llm/chat/completions} desta API, que responde pelo
\cd{AiProvider} em \emph{stream}. Esse \emph{endpoint} não tem pessoa logada, então
leva duas trancas: o \emph{Bearer} \env{AJUDA\_LLM\_SEGREDO} (é o \emph{vendor}) e o
\emph{token} curto da sessão, que a API assinou na abertura e o \emph{vendor} devolve
dentro da mensagem \cd{system} --- o resto daquele \cd{system} é descartado, e o
\emph{prompt} e a base nunca saem da API. E \textbf{o custo não é contado em
\emph{tokens}}: não há \cd{QuotaService}, e sim a reserva de minutos em
\cd{ajuda\_sessoes} contra os tetos de \cd{ajuda\_limites}, feita antes de pedir a
sala. Toda recusa que cabe antes da sala sai em JSON; depois dela, recusa de
negócio vira fala, porque o \emph{vendor} fica mudo diante de qualquer \emph{status}
de erro (medido no \emph{spike}, sem nova tentativa e sem evento ao navegador).

\begin{figure}[H]
\centering
\fitpage{%
\begin{tikzpicture}[x=1cm, y=1cm]

  \tikzset{guarda/.append style={align=left}}
  \def\fundo{-23.55}
  \def\msgp#1#2#3#4#5{\draw[seqmsg] (#1,#3) -- (#2,#3) node[seqlbl, pos=#4, above=0.4pt]{#5};}
  \def\retp#1#2#3#4#5{\draw[seqret] (#1,#3) -- (#2,#3) node[seqlbl, pos=#4, above=0.4pt]{#5};}

  \raia{1.4}{\fundo}{Navegador}{modal da Ajuda}
  \raia{5.4}{\fundo}{\texttt{ajuda/}}{\emph{controllers} + \emph{service}}
  \raiaext{9.0}{\fundo}{PostgreSQL}{Cloud SQL}
  \raiaext{13.4}{\fundo}{LiveAvatar}{HeyGen, FULL Mode}
  \raiaext{17.2}{\fundo}{Gemini}{pelo \texttt{AiProvider}}

  % ---- abrir: tudo o que recusa vem antes da sala
  \msgs{1.4}{5.4}{-1.1}{\texttt{POST /ajuda/sessoes} (BFF)}
  \msgs{5.4}{9.0}{-2.3}{fecha as abertas;\\lê \texttt{ajuda\_limites}\\e \texttt{ajuda\_sessoes}}
  \msgr{9.0}{5.4}{-3.25}{segundos usados no dia\\(fuso da plataforma)}

  \draw[seqframe] (0.4,-3.65) rectangle (6.3,-5.85);
  \node[frameop] at (0.4,-3.65) {break};
  \node[guarda] at (1.55,-4.25) {[desligado, sem atendente,\\mais de 6/min ou menos\\de 1 min no dia]};
  \msgr{5.4}{1.4}{-5.55}{403 $|$ 503 $|$ 429, antes da sala}

  \msgs{5.4}{9.0}{-6.95}{reserva (\texttt{inicio},\\\texttt{duracao\_max\_s})}
  \msgself{5.4}{-7.35}{assina o \emph{token}\\da sessão (HS256)}
  \msgp{5.4}{13.4}{-8.95}{0.7}{\texttt{POST /v1/sessions/token}\\(\texttt{FULL}, com o \emph{token} da sessão)}
  \retp{13.4}{5.4}{-9.65}{0.3}{\texttt{session\_id}, \texttt{session\_token}}
  \msgr{5.4}{1.4}{-10.6}{201 \{\texttt{id}, \texttt{tokenDaSessao},\\\texttt{duracaoMaximaS}, \ldots\}}
  \msgp{1.4}{13.4}{-11.3}{0.85}{\texttt{start()}: entra na sala (WebRTC)}
  \msgp{13.4}{5.4}{-12.0}{0.3}{\texttt{GET /ajuda/llm/models} (\emph{Bearer})}

  % ---- cada fala: o vendor chama o nosso LLM
  \draw[seqframe] (0.4,-12.4) rectangle (18.3,-19.55);
  \node[frameop] at (0.4,-12.4) {loop};
  \node[guarda] at (1.55,-12.8) {[cada fala, por voz\\ou por escrito]};
  \msgp{1.4}{13.4}{-13.5}{0.85}{fala ou \texttt{message(texto)}}
  \msgp{13.4}{5.4}{-14.45}{0.3}{\texttt{POST /ajuda/llm/chat/completions}\\(\emph{Bearer} + \texttt{[sessao:\ldots]} no \texttt{system})}
  \msgself{5.4}{-15.1}{duas trancas\\(senão 401); produto\\ligado? até 20/min?\\(senão vira fala)}
  \msgp{5.4}{17.2}{-16.6}{0.85}{\texttt{conversarEmFluxo}\\(prefixo fixo + papel,\\histórico)}
  \retp{17.2}{5.4}{-17.55}{0.16}{pedaços (SSE do\\Vertex)}
  \retp{5.4}{13.4}{-18.5}{0.7}{\texttt{chat.completion.chunk}\\em SSE, até \texttt{[DONE]}}
  \retp{13.4}{1.4}{-19.2}{0.15}{voz, vídeo e transcrições, pela sala}

  % ---- fechar
  \msgp{1.4}{13.4}{-20.3}{0.85}{\texttt{stop()} (fechar o modal)}
  \msgs{1.4}{5.4}{-21}{\texttt{POST} \ldots\texttt{/encerrar} (\texttt{keepalive})}
  \msgs{5.4}{9.0}{-21.7}{\texttt{fim = agora}}
  \msgp{5.4}{13.4}{-22.4}{0.7}{\texttt{POST /v1/sessions/stop}}

\end{tikzpicture}}
\caption[Sequência: uma sessão da Ajuda com avatar]{Sequência de uma sessão da
  Ajuda com avatar (\adrr{ADR-132}). O navegador fala com a API pelo BFF do hub.
  \texttt{break} é a saída antes da sala, em JSON; dentro do \texttt{loop} só a
  autenticação responde erro (401), e o resto vira fala, porque o \emph{vendor}
  fica mudo diante de qualquer \emph{status} de erro. Fontes:
  \arq{ajuda/ajuda.controller.ts}, \arq{ajuda/ajuda-llm.controller.ts},
  \arq{ajuda/ajuda.service.ts}, \arq{ajuda/regras.ts},
  \arq{ajuda/token-da-sessao.ts}, \arq{ai/liveavatar.client.ts},
  \arq{ai/vertex.provider.ts} e \arq{hub/components/ajuda/use-sessao-da-ajuda.ts}.}
\label{fig:uml-seq-ajuda}
\end{figure}

O que a figura deixa para o texto, para caber legível. Antes do \texttt{break}
rodam as guardas da classe (\cd{OidcGuard}, \cd{RolesGuard} sem o estudante e
\cd{@Produto("ajuda")}). O \emph{token} da sessão carrega
\cd{\{sub, escola, perfil, sid\}} e vai ao \emph{vendor} em
\cd{dynamic\_variables.sessao}. Se o \emph{vendor} não abre a sala, a linha
reservada fecha com duração zero e a resposta é 502 \cd{AVATAR\_FALHOU}. Dentro
do \texttt{loop}, produto desligado ou mais de 20 turnos por minuto viram uma fala
do atendente, em 200, sem chamar o modelo, e a falha do modelo vira a fala de
desculpa. O \emph{vendor} desiste se o primeiro \emph{byte} passar de 10\,s: sem
pedaço em 7\,s, sai a fala de espera e a resposta continua no mesmo turno.
Pergunta e resposta não vão ao log (só a linha \cd{ajuda.turno}, com números). O
encerrar escapa do interruptor (\cd{@SemDisponibilidade}); sessão que fica sem
\cd{fim} conta o teto inteiro, e abrir outra fecha a anterior. O outro atendente
da Ajuda, o mascote (\adrr{ADR-158}), usa a mesma abertura e os mesmos minutos,
mas não tem sala nem \emph{vendor}: recebe um bilhete e conversa por WebSocket
com a própria API (\rota{/ajuda/mascote/ao-vivo}), que fala com o Gemini Live do
Vertex --- é o caminho que a Figura~\ref{fig:uml-implantacao} desenha.

\subsubsection{O caminho de um lote de telemetria}

O diagrama de comunicação mostra os mesmos fatos de um diagrama de sequência,
arrumados de outro jeito: os objetos ficam onde a topologia os põe, as linhas são
os \emph{links} que existem entre eles, e a ordem vai no número de cada mensagem
(\texttt{1.1.2} acontece dentro de \texttt{1.1}, que acontece dentro de \texttt{1}).
Para a telemetria, é a forma certa: o interessante não é o tempo, é
\emph{quem fala com quem}, e três gatilhos independentes (o relógio da SDK, o
\cd{pagehide} e a entrega do Pub/Sub) atravessam os mesmos objetos. A
Figura~\ref{fig:uml-comunicacao} corrige uma leitura comum da tabela de
arquitetura: \textbf{não existe \emph{worker}}. Quem consome o tópico é a própria
API, pela assinatura \emph{push}, e o destino é o PostgreSQL, numa tabela
particionada por mês.

\begin{figure}[H]
\centering
\fitwidth{%
\begin{tikzpicture}[x=1cm, y=1cm]
  \tikzset{ob1/.style={objetoum, text width=34mm, align=center}}

  % grade: colunas x = 1.9 / 9.9 / 16.4; linhas y = 0 / -5 / -10
  \node[ob1] (exp) at (1.9,0) {\underline{\textbf{:Experiência}}\\{\color{unooMuted}código da missão}};
  \node[ob1] (tel) at (1.9,-5.0) {\underline{\textbf{t : TelemetryClient}}\\{\color{unooMuted}\texttt{@edu/sdk}}};
  \node[ob1] (tok) at (9.9,0) {\underline{\textbf{:TokenService}}\\{\color{unooMuted}\texttt{auth/}}};
  \node[ob1] (ctrl) at (9.9,-5.0) {\underline{\textbf{:TelemetryController}}\\{\color{unooMuted}\texttt{telemetry/}}};
  \node[ob1] (ps) at (16.4,-5.0) {\underline{\textbf{:Topic}}\\{\color{unooMuted}\texttt{\{contrato\}-telemetria}}};
  \node[ob1] (guard) at (16.4,-10.0) {\underline{\textbf{:PubSubPushGuard}}\\{\color{unooMuted}ID \emph{token} do Google}};
  \node[ob1, fill=unooGreenWash, draw=unooGreen!60!black] (pg) at (9.9,-10.0)
    {\underline{\textbf{:telemetria\_eventos}}\\{\color{unooMuted}\estereo{TabelaDeTenant}, por mês}};

  % links (linhas retas ou em L; nenhum cruza outro)
  \draw[objetolink] (exp) -- (tel);
  \draw[objetolink] (tel) -- (ctrl);
  \draw[objetolink] (ctrl) -- (tok);
  \draw[objetolink] (ctrl.east |- 0,-4.8) -- (ps.west |- 0,-4.8);
  \draw[objetolink] (ps) -- (guard);
  \draw[objetolink] (guard.west) -- (13.15,-10.0) |- (ctrl.east |- 0,-5.25);
  \draw[objetolink] (ctrl) -- (pg);
  % auto-link do controller (1.1.1 acontece nele)
  \draw[objetolink] ([xshift=-12mm]ctrl.north) -- ++(0,0.7) -- ++(0.75,0) -- ([xshift=-4.5mm]ctrl.north);

  % mensagens numeradas, ao lado do link (o detalhe de cada uma esta na tabela abaixo)
  \msgcom[right]{(2.15,-2.5)}{-90}{\texttt{1: track()}}
  \msgcom[above]{(5.9,-4.78)}{0}{\texttt{1.1: POST /telemetry/batch}}
  \msgcom[below]{(5.9,-5.22)}{0}{\texttt{3: sendBeacon /beacon}}
  \node[msgcomlbl, anchor=east] at (8.25,-3.75) {\texttt{1.1.1: identidade do token}};
  \msgcom[right]{(10.15,-2.5)}{90}{\texttt{3.1: verificar()}}
  \msgcom[above]{(13.15,-4.58)}{0}{\texttt{1.1.2, 3.2:}\\\texttt{publishMessage()}}
  \msgcom[right]{(16.65,-7.5)}{-90}{\texttt{2: push /ingest}}
  \msgcom[right]{(13.4,-7.6)}{90}{\texttt{2.1: ingest()}}
  \msgcom[left]{(9.65,-7.5)}{-90}{\texttt{2.1.1: INSERT ... DO NOTHING}}

  % notas
  \node[actnote, text width=44mm, anchor=north west] at (12.6,0.75)
    {\textbf{Três gatilhos, três números de topo:} \texttt{1} é o relógio da SDK,
     \texttt{3} é a saída da página e \texttt{2} é o \mbox{Pub/Sub}, quando ele quiser.
     Nenhum \mbox{espera} o outro.};
  \node[actnote, text width=40mm, anchor=north west] at (0.0,-8.4)
    {\textbf{Sem \texttt{TELEMETRIA\_TOPIC}} (\emph{compose} local, CI),
     \texttt{1.1.2} vira uma linha de \emph{debug} e nada é gravado.};

\end{tikzpicture}}
\caption[Comunicação: o caminho de um lote de telemetria]{Diagrama de comunicação da
  telemetria. Objetos com nome sublinhado; \emph{links} como linhas; mensagens
  numeradas pela aninhação. O \texttt{ON CONFLICT (event\_id, ts) DO NOTHING}
  torna a reentrega do Pub/Sub inofensiva, e a gravação é agrupada por escola,
  uma transação \cd{comTenancy} por grupo. Fontes:
  \arq{packages/sdk/src/telemetry.ts}, \arq{api/src/telemetry/telemetry.controller.ts},
  \arq{api/src/telemetry/pubsub-push.guard.ts} e \arq{docs/INFRA.md}.}
\label{fig:uml-comunicacao}
\end{figure}

\begin{tabularx}{\linewidth}{@{}L{11mm}L{36mm}X@{}}
\caption{As mensagens da Figura~\ref{fig:uml-comunicacao}}
\label{tab:uml-comunicacao}\\
\toprule
\headrow \thd{Nº} & \thd{Gatilho ou guarda} & \thd{O que acontece}\\
\midrule
\endfirsthead
\toprule
\headrow \thd{Nº} & \thd{Gatilho ou guarda} & \thd{O que acontece}\\
\midrule
\endhead
\texttt{1} & a missão registra um evento & Entra no \emph{buffer}; acima de 500 eventos,
  sai o mais antigo.\\
\zebra \texttt{1.1} & a cada 10 s, ou ao chegar a 50 & \rota{POST /telemetry/batch}
  com \emph{Bearer} e \cd{keepalive}; um lote em voo por vez, e a falha devolve o
  lote ao \emph{buffer} com \emph{backoff} de 1 a 60 s.\\
\texttt{1.1.1} & --- & O Zod descarta \cd{estudanteId} e \cd{escolaId} do corpo; o
  \emph{controller} põe os do \emph{token}. O \cd{estudante\_id} é o \cd{sub}: id
  interno, não nome nem e-mail --- e também não é \emph{hash}.\\
\zebra \texttt{1.1.2} & \env{TELEMETRIA\_TOPIC} definida & \cd{publishMessage(\{json:
  \{eventos\}\})}. A resposta \cd{\{ok, recebidos\}} conta o que foi
  \emph{publicado}, não o que foi gravado.\\
\texttt{2} & entrega do Pub/Sub, \emph{at least once} & \emph{Push} para
  \rota{POST /telemetry/ingest}; o \emph{guard} confere emissor, audiência e o
  e-mail da conta de serviço do \emph{push}. Sem as \emph{envs}, 401.\\
\zebra \texttt{2.1} & o \emph{guard} passou & Decodifica \cd{message.data} (base64) e
  valida com Zod.\\
\texttt{2.1.1} & --- & Agrupa por escola e grava um \cd{INSERT} de várias linhas
  por grupo, dentro de \cd{comTenancy}, com \cd{ON CONFLICT (event\_id, ts) DO
  NOTHING}: reentrega é rotina.\\
\zebra \texttt{3} & \cd{pagehide} ou aba oculta & \cd{sendBeacon} para
  \rota{/telemetry/beacon}, com o \emph{token} no corpo --- o \emph{beacon} não leva
  cabeçalho. Retorno \cd{false} mantém o lote no \emph{buffer}.\\
\texttt{3.1} & --- & \cd{TokenService.verificar} no \emph{token} do corpo; se falha,
  descarta com métrica.\\
\zebra \texttt{3.2} & --- & O mesmo \cd{publishMessage} de \texttt{1.1.2}. A rota
  responde 204 sempre.\\
\bottomrule
\end{tabularx}

\subsubsection{Três relógios: \emph{token}, renovação e sessão}

O diagrama de tempo põe o tempo no eixo horizontal e o estado de cada participante
numa faixa. É o único tipo da UML em que \emph{prazo} é notação de primeira classe
(as restrições entre chaves medem distância no eixo), e é por isso que ele mostra o
que as outras figuras escondem. São três relógios independentes: o \emph{token} de
experiência vive 15 minutos (\arq{hub/lib/chave-token.ts}); a SDK pede um novo a
cada 10, por relógio fixo, sem ler o \cd{exp} do JWT e sem reagir a 401
(\arq{packages/sdk/src/index.ts}); e a sessão do hub morre 2 horas depois da última
renovação (\adrr{ADR-95}), que o próprio pedido de \emph{token} provoca
(\adrr{ADR-96}). A Figura~\ref{fig:uml-tempo} desenha o caso em que o pedido dos
20 minutos falha por rede.

\begin{figure}[H]
\centering
\fitwidth{%
\begin{tikzpicture}[x=1cm, y=1cm]
  % x(t) = 4.3 + 0.32 t, com t em minutos
  \def\T#1{{4.3+0.32*(#1)}}

  % ---- faixas
  \foreach \yt/\yb/\nome in {0/-1.55/Chamadas da\\experiência à API,
                             -1.85/-3.55/\emph{Token} vigente\\na experiência,
                             -3.85/-5.45/Pedido de \emph{token}\\do hub,
                             -5.75/-7.05/Sessão do hub} {
    \draw[tempofaixa] (0,\yt) rectangle (17.6,\yb);
    \node[font=\scriptsize\bfseries\color{unooNavy}, align=left, anchor=west, text width=26mm] at (0.1,{(\yt+\yb)/2}) {\nome};
  }
  \draw[unooAxis, line width=0.6pt] (2.85,0) -- (2.85,-7.05);

  % ---- marcas verticais dos instantes que importam
  \foreach \t in {10,20,25,30} \draw[tempomarca] (\T{\t},0.2) -- (\T{\t},-7.35);

  % ---- faixa 1: 200 ou 401 (estado)
  \node[tempoestado] at (4.15,-0.45) {200};
  \node[tempoestado] at (4.15,-1.15) {401};
  \draw[tempolinha] (\T{0},-0.45) -- (\T{25},-0.45) -- (\T{25},-1.15)
    -- (\T{30},-1.15) -- (\T{30},-0.45) -- (\T{40},-0.45);
  \draw[tempolinha, draw=unooRed] (\T{25},-1.15) -- (\T{30},-1.15);
  \draw[<->, draw=unooRed, line width=0.7pt] (\T{25},-1.38) -- (\T{30},-1.38);
  \node[font=\scriptsize\bfseries\color{unooRed}, anchor=west] at (\T{30.3},-1.15) {\{5 min de 401\}};

  % ---- faixa 2: linha de valor (qual token esta em uso)
  \def\valor#1#2#3{% t0, t1, rotulo
    \draw[tempolinha, line width=0.8pt] (\T{#1+0.35},-2.2) -- (\T{#2-0.35},-2.2)
      -- (\T{#2},-2.525) -- (\T{#2-0.35},-2.85) -- (\T{#1+0.35},-2.85) -- (\T{#1},-2.525) -- cycle;
    \node[font=\scriptsize] at (\T{(#1+#2)/2},-2.525) {#3};}
  \valor{0}{10}{\texttt{A}: vale até 15}
  \valor{10}{30}{\texttt{B}: vale até 25}
  \valor{30}{40}{\texttt{D}: vale até 45}
  \fill[pattern={Lines[angle=45, distance=3pt, line width=0.4pt]}, pattern color=unooRed]
    (\T{25},-2.2) rectangle (\T{30-0.35},-2.85);
  \draw[<->, draw=unooNavy, line width=0.6pt] (\T{10},-3.3) -- (\T{25},-3.3)
    node[restricao, midway]{\{15 min\}};

  % ---- faixa 3: o pedido de token (pulsos a cada 10 min)
  \node[tempoestado] at (4.15,-4.2) {ocioso};
  \node[tempoestado] at (4.15,-4.95) {pedindo};
  \draw[tempolinha] (\T{0},-4.2)
    -- (\T{10},-4.2) -- (\T{10},-4.95) -- (\T{10.6},-4.95) -- (\T{10.6},-4.2)
    -- (\T{20},-4.2) -- (\T{20},-4.95) -- (\T{20.6},-4.95) -- (\T{20.6},-4.2)
    -- (\T{30},-4.2) -- (\T{30},-4.95) -- (\T{30.6},-4.95) -- (\T{30.6},-4.2)
    -- (\T{40},-4.2) -- (\T{40},-4.95) -- (\T{40.4},-4.95);
  \node[font=\scriptsize\color{unooGreen!70!black}, anchor=south west] at (\T{10.7},-4.95) {ok};
  \node[font=\scriptsize\bfseries\color{unooRed}, anchor=south west] at (\T{20.7},-4.95) {falhou (rede)};
  \node[font=\scriptsize\color{unooGreen!70!black}, anchor=south west] at (\T{30.7},-4.95) {ok};
  \draw[<->, draw=unooNavy, line width=0.6pt] (\T{10},-5.2) -- (\T{20},-5.2)
    node[restricao, midway]{\{10 min\}};

  % ---- faixa 4: a sessao desliza a cada pedido que chega
  \node[tempoestado] at (4.15,-6.1) {ativa};
  \node[tempoestado] at (4.15,-6.7) {expirada};
  \draw[tempolinha] (\T{0},-6.1) -- (\T{40},-6.1);
  \foreach \t in {0,10,30} \fill[unooNavy] (\T{\t},-6.1) circle (1.6pt);
  \node[font=\scriptsize\color{unooInkSoft}, anchor=north west, align=left] at (\T{0.3},-6.2)
    {cada ponto renova o prazo: \{expira 2 h depois do último\}};

  % ---- mensagens entre faixas (os pedidos e as respostas)
  \foreach \t in {10,20,30} \draw[seqmsg, line width=0.6pt] (\T{\t-0.18},-0.45) -- (\T{\t-0.18},-4.2);
  \node[font=\scriptsize, anchor=east, fill=unooSurface, inner sep=1pt] at (\T{9.75},-0.82) {\texttt{exp:token-expirando}};
  \foreach \t in {10.6,30.6} \draw[seqret, line width=0.6pt] (\T{\t},-4.2) -- (\T{\t},-2.85);
  \node[font=\scriptsize, anchor=west, fill=unooSurface, inner sep=1pt] at (\T{10.75},-3.03) {\texttt{hub:token}};

  % ---- eixo do tempo
  \draw[->, draw=unooInkSoft, line width=0.7pt] (4.1,-7.4) -- (17.5,-7.4);
  \foreach \t in {0,5,...,40} {
    \draw[unooInkSoft] (\T{\t},-7.34) -- (\T{\t},-7.46);
    \node[tempotick, below] at (\T{\t},-7.44) {\t};
  }
  \node[tempotick, anchor=west] at (17.5,-7.4) {min};

\end{tikzpicture}}
\caption[Tempo: \emph{token}, renovação e sessão]{Diagrama de tempo de uma missão
  aberta, no cenário em que o pedido de \emph{token} dos 20 minutos falha por rede.
  O \emph{token} \texttt{B} (emitido aos 10) vence aos 25; o próximo pedido é aos
  30; entre os dois, toda chamada da experiência à API recebe 401. As restrições
  entre chaves são de duração. Fontes: \arq{packages/sdk/src/index.ts},
  \arq{hub/components/ExpFrame.tsx}, \arq{hub/lib/chave-token.ts},
  \arq{hub/app/api/exp-token/route.ts} e \arq{hub/lib/sessao.ts}.}
\label{fig:uml-tempo}
\end{figure}

\begin{limitacao}
\textbf{A falha transitória da renovação custa 5 minutos de 401, não zero.} O
comentário do \cd{ExpFrame} diz que \enquote{o \emph{token} atual ainda vale alguns
minutos e a SDK pede de novo em 10} --- as duas metades são verdade, mas os
números não fecham: sobram 5 minutos de validade para 10 de espera. O efeito na
sala é uma missão que falha por 5 minutos e volta sozinha, sem aviso, e a
experiência vê só \enquote{API \ldots{} \seta{} 401}. Consertos possíveis, todos
fora desta frente: a SDK reagir a 401 pedindo \emph{token} na hora, o
\cd{ExpFrame} tentar de novo em segundos em vez de esperar o próximo tique, ou o
relógio da SDK ler o \cd{exp} do JWT. O inverso também vale e é a consequência
aceita na \adrr{ADR-96}: enquanto a aba estiver aberta e a rede responder, a
sessão nunca expira, porque cada tique a renova.
\end{limitacao}

\subsubsection{A ponte como máquina de estados de protocolo}

A sequência do \emph{handshake} (Figura~\ref{fig:seq-handshake}) mostra \emph{uma}
conversa que dá certo. A máquina de estados de protocolo da
Figura~\ref{fig:uml-estados-ponte} mostra \emph{todas} as conversas permitidas,
do lado da experiência: em que estado a SDK pode estar, que mensagem é aceita em
cada um e o que acontece com o resto. Não há \cd{enum} de estado no código ---
\arq{packages/sdk/src/bridge.ts} guarda a máquina em uma \emph{flag}
(\cd{conectado}) e dois temporizadores ---, e é por isso que desenhá-la vale: é o
único lugar em que os quatro estados aparecem juntos.

\begin{figure}[H]
\centering
\fitwidth{%
\begin{tikzpicture}[x=1cm, y=1cm]
  \tikzset{
    estc/.style={rectangle split, rectangle split parts=2, rounded corners=6pt,
      draw=unooBlue, line width=0.85pt, fill=unooBlueWash, font=\scriptsize,
      text width=58mm, inner sep=4pt,
      every text node part/.style={align=center},
      every two node part/.style={align=left}},
    esterr/.style={estado, text width=34mm, fill=unooRedWash, draw=unooRed},
  }

  % ---- o estado inicial
  \node[estc, anchor=north west] (agu) at (0.2,-0.6)
    {\textbf{AguardandoInit}
     \nodepart{two}
     \textit{entry} / \texttt{normalizarHubOrigin}; envia \texttt{exp:ready}\\
     \textit{do} / reenvia \texttt{exp:ready} a cada 400 ms\\
     \texttt{message} [\texttt{origin} $\neq$ \texttt{hubOrigin}] / descarta\\
     \texttt{message} [não casa \texttt{HubInitSchema}] / ignora};
  \node[actstart] (ini) at ([xshift=9mm, yshift=11mm]agu.north west) {};

  % ---- conectada: estado composto
  \draw[estadocomposto] (9.0,0.9) rectangle (17.8,-7.6);
  \node[font=\scriptsize\bfseries, anchor=north west] at (9.15,0.8) {Conectada};
  \node[font=\scriptsize, anchor=north west, text width=84mm, align=left] at (9.15,0.35)
    {\texttt{hub:token} / \texttt{token := msg.token}\\
     \textit{after}(10 min) / envia \texttt{exp:token-expirando}\\
     \texttt{hub:init} / ignora (resposta a um \texttt{exp:ready} reenviado)\\
     \texttt{message} [\texttt{origin} $\neq$ \texttt{hubOrigin}] / descarta\\
     \texttt{message} [não casa \texttt{HubParaExpSchema}] / avisa e descarta};
  \draw[draw=unooBlue, line width=0.6pt] (9.0,-1.85) -- (17.8,-1.85);
  \node[actstart, minimum size=3.6mm] (ini2) at (9.5,-2.7) {};
  \node[estado, text width=26mm] (ativa) at (11.55,-4.0) {Ativa};
  \node[estado, text width=26mm] (pausa) at (16.15,-4.0) {Pausada};
  \draw[trans] (ini2) |- (ativa.west);
  \draw[trans] ([yshift=1.5mm]ativa.east) -- ([yshift=1.5mm]pausa.west)
    node[translbl, midway, above]{\texttt{hub:pause}};
  \draw[trans] ([yshift=-1.5mm]pausa.west) -- ([yshift=-1.5mm]ativa.east)
    node[translbl, midway, below]{\texttt{hub:resume}};
  \node[font=\scriptsize\color{unooInkSoft}, anchor=north west, text width=82mm, align=flush left] at (9.2,-5.2)
    {A SDK não guarda \emph{Ativa} nem \emph{Pausada}: repassa as duas mensagens aos
     \emph{handlers} \texttt{onPause}/\texttt{onResume}, e a experiência para os
     próprios cronômetros. O hub manda \texttt{hub:pause} em
     \texttt{visibilitychange}, em \texttt{blur} (salvo quando o foco foi para o
     próprio \texttt{iframe}) e antes do aviso de sessão expirada.};

  % ---- os dois fins com erro
  \node[esterr, anchor=north] (inc) at (1.95,-5.0) {\textbf{Incompatível}\\{\color{unooInkSoft}\texttt{reject}:\\\enquote{Protocolo incompatível}}};
  \node[esterr, anchor=north] (esg) at (6.0,-5.0) {\textbf{TempoEsgotado}\\{\color{unooInkSoft}\texttt{reject}:\\\enquote{Timeout no \emph{handshake}}}};
  \node[actend] (f1) at ([yshift=-11mm]inc.south) {};
  \node[actend] (f2) at ([yshift=-11mm]esg.south) {};

  % ---- transicoes
  \draw[trans] (ini) -- (ini |- agu.north)
    node[translbl, midway, right=2pt, align=left]{\texttt{connectToHub}(\texttt{hubOrigin})};
  \draw[trans] (agu.east |- 0,-1.6) -- (9.0,-1.6)
    node[translbl, midway, above, align=center]{\texttt{hub:init}\\{[}\texttt{trunc(v) = 2}{]}}
    node[translbl, midway, below, align=center]{/ limpa os\\temporizadores;\\\texttt{resolve}};
  \draw[trans] (inc.north |- agu.south) -- (inc.north)
    node[translbl, midway, right=1pt, align=left]{\texttt{hub:init}\\{[}\texttt{trunc(v)} $\neq$ \texttt{2}{]}};
  \draw[trans] (esg.north |- agu.south) -- (esg.north)
    node[translbl, midway, right=1pt, align=left]{\textit{after}(8 s)};
  \draw[trans] (inc) -- (f1);
  \draw[trans] (esg) -- (f2);

  \node[actnote, text width=84mm, align=flush left, anchor=north west] at (9.0,-8.0)
    {\textbf{Na prática, o incompatível termina em tempo esgotado.} O
     \texttt{ExpFrame} confere a versão do \texttt{exp:ready} e, se o maior
     diverge, registra \texttt{protocolo-incompativel} e \emph{não responde}. A
     SDK nunca recebe o \texttt{hub:init} que a levaria a \emph{Incompatível}: cai
     no \emph{after}(8 s). O estado existe para o dia em que um hub responder com
     outra versão.};

\end{tikzpicture}}
\caption[Máquina de estados de protocolo: a ponte]{Máquina de estados de protocolo da
  ponte, do lado da experiência (\cd{connectToHub} e \cd{createExperiencia}). Os
  compartimentos listam as transições internas --- as que não mudam de estado ---
  no formato \emph{evento} [\emph{guarda}] / \emph{efeito}. Fontes:
  \arq{packages/sdk/src/bridge.ts}, \arq{packages/sdk/src/index.ts},
  \arq{packages/sdk/src/types.ts} (\cd{PROTOCOL\_VERSION = 2}) e
  \arq{hub/components/ExpFrame.tsx}.}
\label{fig:uml-estados-ponte}
\end{figure}

\subsubsection{A jornada do estudante, do \emph{login} à missão}

O diagrama de visão geral de interação é um diagrama de atividades em que cada nó
é uma \emph{interação}: um quadro \texttt{ref} aponta para uma sequência desenhada
em outro lugar, e o fluxo entre eles tem as decisões, as bifurcações e os fins de
um diagrama de atividades. É o índice das interações das seções anteriores. A
Figura~\ref{fig:uml-visao-geral} costura a jornada de um estudante numa manhã de
aula: cada quadro diz em que figura a conversa está desenhada por inteiro.

\begin{figure}[H]
\centering
\fitpage{%
\begin{tikzpicture}[x=1cm, y=1cm]
  \tikzset{
    rf/.style={refint, text width=#1, inner ysep=4.5pt,
      label={[frameop, anchor=north west, inner sep=1.4pt]north west:ref}},
    rf/.default=44mm,
    dc/.style={decision, text width=24mm, aspect=2.4},
    flowfinal/.style={circle, draw=unooInk, line width=0.9pt, fill=unooSurface,
      minimum size=4.6mm, inner sep=0pt,
      path picture={\draw[unooInk, line width=0.8pt]
        (path picture bounding box.south west) -- (path picture bounding box.north east)
        (path picture bounding box.north west) -- (path picture bounding box.south east);}},
  }

  \node[actstart] (ini) at (8.2,0.4) {};
  \node[rf] (login) at (8.2,-0.9)
    {\emph{login} (\ucr{UC-01})\\\texttt{signIn("credentials")} \seta{} \texttt{POST /auth/login}};
  \node[dc] (dsenha) at (8.2,-2.55) {senha provisória?};
  \node[rf=36mm] (senha) at (13.6,-2.55) {trocar a senha\\\texttt{/trocar-senha} (\ucr{UC-25})};
  \node[escolha] (m1) at (8.2,-3.85) {};
  \node[rf] (grade) at (8.2,-5.0)
    {a grade do estudante\\\texttt{GET /salas/minhas}, sem as indisponíveis};
  \node[dc] (dlib) at (8.2,-6.65) {o \emph{card} está liberado?};
  \node[action, text width=34mm] (cad) at (2.6,-6.65) {cadeado: \enquote{Ainda não abriu} ou \enquote{Esta você já fechou}};
  \node[flowfinal] (ff1) at (2.6,-7.85) {};
  \node[rf] (entra) at (8.2,-8.55)
    {entrada na missão\\as duas travas (Figura~\ref{fig:uml-seq-entrada})};
  \node[rf] (hand) at (8.2,-9.95)
    {\emph{handshake}\\\texttt{exp:ready}/\texttt{hub:init} (Figura~\ref{fig:seq-handshake})};
  \node[bar, minimum width=120mm] (fork) at (8.2,-11.05) {};
  \node[rf=40mm] (dom) at (2.9,-12.3)
    {rotas de domínio\\e IA com quota (Fig.~\ref{fig:uml-seq-ia})};
  \node[rf=40mm] (tel) at (8.2,-12.3)
    {telemetria em lote\\(Figura~\ref{fig:uml-comunicacao})};
  \node[rf=40mm] (ren) at (13.5,-12.3)
    {renovação do \emph{token}\\a cada 10 min (Fig.~\ref{fig:uml-tempo})};
  \node[bar, minimum width=120mm] (join) at (8.2,-13.55) {};
  \node[dc] (dfim) at (8.2,-14.8) {como a missão\\termina?};
  \node[rf=40mm] (conc) at (2.9,-17.2)
    {concluir\\\texttt{POST .../progresso/}\\\texttt{:expId/concluir}};
  \node[actionerr, text width=38mm] (exp) at (8.2,-17.2) {aviso \enquote{Sua sessão expirou};\\\texttt{hub:pause}};
  \node[actionerr, text width=38mm] (trav) at (13.5,-17.2) {\enquote{A missão travou},\\com o \emph{hash} do erro};
  \node[actend] (f1) at (2.9,-18.6) {};
  \node[actend] (f2) at (8.2,-18.6) {};
  \node[actend] (f3) at (13.5,-18.6) {};

  \draw[flow] (ini) -- (login);
  \draw[flow] (login) -- (dsenha);
  \draw[flow] (dsenha) -- (senha) node[flowlbl, midway, above]{[sim]};
  \draw[flow] (senha.south) |- (m1.east);
  \draw[flow] (dsenha) -- (m1) node[flowlbl, midway, left=2pt]{[não]};
  \draw[flow] (m1) -- (grade);
  \draw[flow] (grade) -- (dlib);
  \draw[flow] (dlib) -- (cad) node[flowlbl, midway, above]{[não]};
  \draw[flow] (cad) -- (ff1);
  \draw[flow] (dlib) -- (entra) node[flowlbl, midway, left=2pt]{[sim]};
  \draw[flow] (entra) -- (hand);
  \draw[flow] (hand) -- (fork);
  \foreach \n in {dom, tel, ren} {
    \draw[flow] (fork.south -| \n.north) -- (\n.north);
    \draw[flow] (\n.south) -- (join.north -| \n.south);
  }
  \draw[flow] (join) -- (dfim);
  \draw[flow] (dfim.west) -| (conc.north) node[flowlbl, pos=0.25, above]{[\texttt{exp:concluir}]};
  \draw[flow] (dfim.south) -- (exp.north) node[flowlbl, pos=0.55, right=2pt, align=left]{[\texttt{/api/exp-token}\\dá 401 ou 403]};
  \draw[flow] (dfim.east) -| (trav.north) node[flowlbl, pos=0.25, above]{[\texttt{exp:erro}]};
  \draw[flow] (conc) -- (f1);
  \draw[flow] (exp) -- (f2);
  \draw[flow] (trav) -- (f3);

\end{tikzpicture}}
\caption[Visão geral de interação: a jornada do estudante]{Visão geral de interação da
  jornada do estudante. Cada quadro \texttt{ref} é um uso de interação desenhada
  por inteiro na figura indicada; as barras são bifurcação e junção --- as três
  interações do meio correm juntas enquanto a missão está aberta. A saída pelo
  cadeado termina em \emph{X} (fim de fluxo): o estudante continua na \emph{home}.
  Fontes: \arq{hub/app/(plataforma)/login/login-form.tsx},
  \arq{hub/middleware.ts}, \arq{hub/components/cartao-missao-estudante.tsx},
  \arq{hub/components/ExpHost.tsx} e \arq{hub/components/ExpFrame.tsx}.}
\label{fig:uml-visao-geral}
\end{figure}

% ---------------------------------------------------------------------------
\subsection{Visão física: a implantação de um contrato}
\label{sec:uml-fisica}

O Capítulo~\ref{cap:arqnuvem} desenha a implantação com a notação de
infraestrutura: rede, sub-rede, serviço gerenciado, a seta pública e a privada. O
diagrama de implantação da UML (Figura~\ref{fig:uml-implantacao}) responde outra
pergunta: \emph{que artefato roda em que nó}. Um nó é um recurso computacional
(caixa em perspectiva); um ambiente de execução é um nó que hospeda software;
um \textbf{artefato} é o arquivo que o CI produz e alguém executa, e é ele que
liga o diagrama ao \emph{pipeline} do Capítulo~\ref{cap:entrega}. Três leituras
saem daqui. A imagem da API é implantada \emph{duas} vezes: no serviço e no
\emph{job} de migração, que roda \cd{prisma migrate deploy} na VPC antes de todo
\emph{deploy} da API. O \emph{bundle} é armazenado num nó e executado em outro: vive
no \emph{bucket}, roda no navegador, e a versão que vale não está no \emph{bucket},
está na linha de \cd{experiencias} na API. E nenhum artefato é implantado à mão:
cada seta \estereo{implanta} que sai do GitHub Actions é um passo do \arq{deploy.yml} --- inclusive a do
agente de IA do professor, cuja imagem não roda num Cloud Run, e sim no Vertex AI
Agent Engine, chamado só pela API. O único nó externo que recebe tráfego
\emph{e} chama a API é o do LiveAvatar, da Ajuda com avatar (\adrr{ADR-132}): nada
dele é implantado pelo CI --- o contexto e a configuração de LLM do \emph{vendor}
saem de \arq{scripts/liveavatar-configurar.mjs}, um conjunto por contrato e
ambiente.

\begin{figure}[H]
\centering
\fitpage{%
\begin{tikzpicture}[x=1cm, y=1cm]
  \tikzset{
    art/.style={artefato, text width=#1},
    art/.default=28mm,
    nl/.style={font=\scriptsize, anchor=north west, align=left, text width=#1, inner sep=2pt},
    nl/.default=31mm,
    ext/.style={fill=unooPaper, draw=unooInkSoft},
  }
  % grade: coluna externa x 0-3.4; dentro da nuvem B 5.2-8.6, C 10.2-13.6, D 14.9-18.3
  %        linhas R1 0/-2.6, R2 -3.6/-7.4, R3 -8.4/-10.8, R4 -11.8/-14.2, R5 -15.2/-16.8

  % ---- fora do Google: o vendor da Ajuda e o navegador
  \nodo[ext]{lav}{(0,2.5)}{(3.4,0.7)}{\estereo{externo} \textbf{LiveAvatar}\\HeyGen, na AWS}
  \nodo{nav}{(0,-0.2)}{(3.4,-7.4)}{\estereo{device}\\\textbf{Navegador} da escola}
  \nodo[fill=unooBlueWash]{navenv}{(0.2,-1.2)}{(3.1,-3.3)}{\estereo{browser} aba do hub}
  \node[art=22mm] at (1.6,-2.5) {\estereo{artefato}\\página do hub};
  \nodo[fill=unooOrangeWash, draw=unooOrange]{ifr}{(0.2,-3.9)}{(3.1,-7.15)}{\estereo{sandbox} \texttt{iframe}}
  \node[font=\scriptsize\color{unooMuted}, align=center] at (1.6,-4.85) {\emph{token} só\\em memória};
  \node[art=22mm] (bnav) at (1.6,-6.3) {\estereo{Bundle}\\\texttt{\{expId\}/}\\\texttt{\{versao\}/}};
  \nodo[ext]{gha}{(0,-11.8)}{(3.4,-14.2)}{\estereo{nó} \textbf{GitHub}\\\textbf{Actions}}
  \node[nl=30mm] at (0.05,-12.75) {\texttt{deploy.yml}; WIF,\\sem chave JSON};

  % ---- a stack do contrato
  \draw[nuvem] (4.6,0.8) rectangle (19.0,-17.3);
  \node[nuvemlbl] at (4.8,0.95) {\estereo{projeto GCP} uma \emph{stack} por contrato};

  % R1
  \nodo{lb}{(5.2,0)}{(8.6,-2.6)}{\estereo{nó} \textbf{Load Balancer}\\global + Cloud CDN}
  \node[nl] at (5.25,-1.25) {um IP; \emph{host rules}:\\\texttt{hub}, \texttt{api}, \texttt{exp}, \texttt{assets}};
  \nodo[fill=unooBlueWash]{runhub}{(10.2,0)}{(13.6,-2.6)}{\estereo{Cloud Run}\\\texttt{\{contrato\}-frontend}}
  \node[art] at (11.9,-1.75) {\estereo{imagem}\\\texttt{hub:\{sha\}}\\\texttt{apps/hub/server.js}};
  \nodo[fill=unooBlueWash]{job}{(14.9,0)}{(18.3,-2.6)}{\estereo{Cloud Run Job}\\\texttt{\{contrato\}-migrate}}
  \node[art] at (16.6,-1.75) {\estereo{imagem}\\\texttt{api:\{sha\}}\\\texttt{prisma migrate deploy}};

  % R2
  \nodo{gcs}{(5.2,-3.6)}{(8.6,-7.4)}{\estereo{nó} \textbf{Cloud Storage}}
  \node[nl] at (5.25,-4.1) {\texttt{-exp}, \texttt{-assets}: públicos;\\\texttt{-planos} e mapas:\\\textbf{privados}};
  \node[art] (bgcs) at (6.9,-6.3) {\estereo{Bundle}\\\texttt{\{expId\}/\{versao\}/}\\{\color{unooMuted}\texttt{index.html}:}\\{\color{unooMuted}\texttt{max-age=60}}};
  \nodo[fill=unooBlueWash]{runapi}{(10.2,-3.6)}{(13.6,-7.4)}{\estereo{Cloud Run}\\\texttt{\{contrato\}-api}}
  \node[art] at (11.9,-5.75) {\estereo{imagem}\\\texttt{api:\{sha\}}\\\texttt{dist/src/main.js}};
  \nodo{sql}{(14.9,-3.6)}{(18.3,-7.4)}{\estereo{nó} \textbf{Cloud SQL}\\\texttt{\{contrato\}-pg}}
  \node[nl] at (14.95,-4.45) {PostgreSQL, HA + PITR;\\só IP privado; \emph{role}\\\texttt{app} sem \texttt{BYPASSRLS};\\a posse das conversas\\do agente};

  % R3
  \nodo{ps}{(5.2,-8.4)}{(8.6,-10.8)}{\estereo{nó} \textbf{Pub/Sub}}
  \node[nl] at (5.25,-8.85) {\emph{topic} \texttt{-telemetria};\\\emph{push}\\\texttt{-telemetria-ingest}};
  \nodo[ext]{vx}{(10.2,-8.4)}{(13.6,-10.8)}{\estereo{serviço gerenciado}\\\textbf{Vertex AI}}
  \node[nl] at (10.25,-9.25) {Gemini em \texttt{global}\\(reserva \texttt{us}); Live\\em \texttt{us-central1}};
  \nodo[fill=unooBlueWash]{eng}{(14.9,-8.4)}{(18.3,-10.8)}{\estereo{Agent Engine} \texttt{us-east1}}
  \node[art=30mm] (aeng) at (16.6,-9.45) {\estereo{imagem}\\\texttt{agente-professor:}\\\texttt{\{sha\}}};
  \node[font=\scriptsize] at (16.6,-10.5) {SA \texttt{\{prefixo\}-agente}};

  % R4
  \nodo{ar}{(5.2,-11.8)}{(8.6,-14.2)}{\estereo{nó} \textbf{Artifact Registry}}
  \node[nl] at (5.25,-12.25) {\texttt{hub}, \texttt{api} e\\\texttt{agente-professor},\\por \texttt{\{sha\}}};
  \nodo{bag}{(10.2,-11.8)}{(13.6,-14.2)}{\estereo{nó} \textbf{Cloud Storage}\\\texttt{-agente-professor}}
  \node[nl] at (10.25,-12.65) {privado: BNCC, livros\\e habilidades};
  \nodo[ext]{vs}{(14.9,-11.8)}{(18.3,-14.2)}{\estereo{serviço gerenciado}\\\textbf{Vertex AI Search}}
  \node[nl] at (14.95,-12.6) {\emph{data stores} em\\\texttt{global}: \texttt{-habilidades}\\(estruturado) e\\\texttt{-acervo} (PDFs)};

  % R5
  \nodo{sm}{(5.2,-15.2)}{(8.6,-16.8)}{\estereo{nó} \textbf{Secret Manager}}
  \node[nl] at (5.25,-15.6) {dez segredos, \emph{env}\\do hub e da API};

  % ---- caminhos de comunicacao
  \draw[caminhocom] (lav.south) -- (lav.south |- nav.north)
    node[cxlabel, midway, right=2pt]{WebRTC (vídeo)};
  \draw[caminhocom] (nav.east |- 0,-1.3) -- (lb.west |- 0,-1.3)
    node[cxlabel, pos=0.3, above=1pt, align=center]{HTTPS\\e WSS};
  \draw[caminhocom] (lb.east |- 0,-1.0) -- (runhub.west |- 0,-1.0)
    node[cxlabel, midway, above]{\texttt{hub}};
  \draw[caminhocom] (lb.east |- 0,-2.1) -- (9.3,-2.1) -- node[cxlabel, pos=0.45, right=1pt]{\texttt{api}} (9.3,-4.5) -- (runapi.west |- 0,-4.5);
  \draw[caminhocom] (lb.south) -- (gcs.north) node[cxlabel, midway, right=2pt]{CDN};
  \draw[caminhocom] (12.9,-2.6) -- (12.9,-3.6)
    node[cxlabel, midway, left=2pt]{\texttt{API\_INTERNAL\_URL}};
  \draw[caminhocom] (runapi.west |- 0,-5.6) -- (gcs.east |- 0,-5.6)
    node[cxlabel, midway, above=1pt, align=center]{\emph{assets},\\planos};
  \draw[caminhocom] (runapi.west |- 0,-6.7) -- (9.6,-6.7) -- node[cxlabel, pos=0.5, align=center]{publica;\\\emph{push}} (9.6,-9.6) -- (ps.east |- 0,-9.6);
  \draw[caminhocom] (runapi.east |- 0,-5.6) -- (sql.west |- 0,-5.6) node[cxlabel, midway, above]{VPC};
  \draw[caminhocom] (job.south) -- (sql.north) node[cxlabel, midway, right=2pt]{VPC};
  \draw[caminhocom] (11.0,-7.4) -- (11.0,-8.4) node[cxlabel, pos=0.45, right=2pt]{Gemini e Live};
  \draw[caminhocom] (13.1,-7.4) -- (13.1,-7.85) -- node[cxlabel, pos=0.5]{SSE (\texttt{:streamQuery})} (16.6,-7.85) -- (16.6,-8.4);
  \draw[caminhocom] (eng.west |- 0,-10.2) -- (vx.east |- 0,-10.2) node[cxlabel, midway, above]{Gemini};
  \draw[caminhocom] (eng.south) -- (vs.north) node[cxlabel, midway, right=2pt]{busca};
  \draw[caminhocom] (vs.west |- 0,-13.3) -- (bag.east |- 0,-13.3) node[cxlabel, midway, above]{lê};
  \draw[caminhocom] (lav.east |- 0,1.6) -- node[cxlabel, pos=0.5, above=1pt]{HTTPS: \texttt{/v1/sessions} (API \seta{} \emph{vendor}) e \texttt{/ajuda/llm} (\emph{vendor} \seta{} LB \seta{} API)} (14.35,1.6) -- (14.35,-4.4) -- (runapi.east |- 0,-4.4);

  % ---- o CI implanta cada artefato
  \draw[implanta] (bgcs.west) -- (bnav.east) node[cxlabel, pos=0.68, above=1pt]{baixa};
  \draw[implanta] (gha.east |- 0,-12.3) -- (4.1,-12.3) -- (4.1,-7.0) -- (gcs.west |- 0,-7.0);
  \node[cxlabel, anchor=east] at (4.0,-9.6) {\texttt{rsync}};
  \draw[implanta] (gha.east |- 0,-13.4) -- (ar.west |- 0,-13.4) node[cxlabel, pos=0.35, above=1pt]{\emph{push}};
  \draw[implanta] (gha.south) -- (gha.south |- 0,-17.8) -- node[cxlabel, pos=0.5, above=1pt]{\texttt{PATCH} da \texttt{imageUri}} (18.75,-17.8) -- (18.75,-10.2) -- (eng.east |- 0,-10.2);

\end{tikzpicture}}
\begin{legendadiag}
\mbox{\amostra{nodocorpo} nó}\quad
\mbox{\amostra{nodocorpo, fill=unooBlueWash} ambiente de execução}\quad
\mbox{\amostra{nodocorpo, fill=unooPaper, draw=unooInkSoft} externo ou gerenciado}\quad
\mbox{\amostra{nodocorpo, fill=unooOrangeWash, draw=unooOrange} \emph{sandbox}}\quad
\mbox{\amostra{artefato} artefato}\quad
\mbox{\tikz[baseline=-0.55ex]\draw[caminhocom] (0,0) -- (0.6,0); caminho de comunicação}\quad
\mbox{\tikz[baseline=-0.55ex]\draw[implanta] (0,0) -- (0.6,0); \estereo{implanta}}
\end{legendadiag}
\caption[Implantação: que artefato roda em que nó]{Diagrama de implantação de um
  contrato. A mesma imagem \texttt{api:\{sha\}} está implantada em dois nós, e o
  \emph{bundle} está armazenado num e executado em outro. O CI ainda executa o
  \emph{job} de migração e publica o manifesto (\rota{PUT /experiencias/:expId}),
  o que o diagrama deixa no texto para não cruzar a figura. Escala, rede e
  certificados são do Terraform, em outro repositório
  (Capítulo~\ref{cap:arqnuvem}). Fontes: \arq{apps/api/Dockerfile},
  \arq{apps/hub/Dockerfile}, \arq{.github/workflows/deploy.yml},
  \arq{docs/regras/agente-professor.md}, \arq{scripts/liveavatar-configurar.mjs}
  e \arq{docs/INFRA.md}.}
\label{fig:uml-implantacao}
\end{figure}

A parte de IA da figura tem três regras que o desenho só sugere. \textbf{O agente
de IA do professor só é alcançado pela API}: a imagem \texttt{agente-professor}
roda no Vertex AI Agent Engine (\texttt{us-east1} em dev e hml) com a conta de
serviço \cd{\{prefixo\}-agente}, e a API o chama com o ADC da conta dela
(\cd{:query} e \cd{:streamQuery?alt=sse}); o navegador nunca fala com ele nem com
os \emph{data stores}. O agente usa o Gemini em \cd{global}, com reserva em
\cd{us}, e os dois \emph{data stores} do Vertex AI Search em \cd{global}
(\cd{agente-professor-habilidades}, estruturado, e \cd{agente-professor-acervo},
os PDFs), que leem o \emph{bucket} privado \cd{<projeto>-agente-professor}; a
posse de cada conversa mora no PostgreSQL (\cd{agente\_professor\_conversas}). O
nó do Agent Engine só existe depois do segundo \emph{apply} do Terraform; daí em
diante o CI troca a imagem dele por \cd{PATCH} da \cd{imageUri}
(\adrr{ADR-126}). \textbf{A Ajuda tem dois atendentes e nenhum recurso próprio}:
o LiveAvatar (\adrr{ADR-132}) fica fora do Google e conversa com a \emph{stack}
pelos dois lados --- a API pede e fecha a sala em \cd{api.liveavatar.com}, e o
\emph{vendor} chama \rota{/ajuda/llm} pela URL pública da API, o mesmo LB ---,
com o vídeo direto ao navegador; o mascote (\adrr{ADR-158}) é o WSS do navegador à
API pelo mesmo LB, e da API ao Gemini Live do Vertex em \cd{us-central1}. E
\textbf{o WSO2 não aparece}: é o emissor futuro, num projeto próprio; hoje quem
assina os \emph{tokens} é o hub.
